% ============================================================
%  3. ESQUEMA DE SOLUCIÓN Y ALCANCE  (Subdocumento 3 — Formulario T-12)
%  Informe 1 — Esquema de solución y del alcance
%  Fuente: Informe (1).md (Solución propuesta y Alcances)
% ============================================================
\chapter{Esquema de Solución y Alcance}
\label{sec:solucion}

\section{Lo que haremos en el proyecto}

La solución propuesta para la Distribuidora Puelche S.A. consiste en la modernización
integral de la operación logística de la cadena de frío, la trazabilidad sanitaria y la
gestión de preventa y reparto, sobre una plataforma digital de misión crítica de
carácter obligatoriamente híbrido. La implantación se estructura en dos etapas alineadas
con las restricciones operacionales y de cronograma del caso.

La estrategia de la solución se funda en cinco ejes coherentes con el problema
diagnosticado:

\begin{enumerate}
  \item \textbf{Operación desconectada (\emph{offline-first}) como regla de diseño:}
        ante la ausencia de cobertura en rutas rurales, cámaras de congelado y cortes de
        enlace, toda capacidad crítica opera localmente y sincroniza de forma
        determinista al recuperar conectividad.
  \item \textbf{Trazabilidad lote-a-lote desde el origen:} el registro estructurado del
        lote en el punto de recepción, el monitoreo continuo de temperatura y la
        consulta de clientes afectados habilitan la respuesta a retiros sanitarios en
        horas, no en días.
  \item \textbf{Preventa y reparto con información en el punto de decisión:} stock,
        crédito, deuda vencida e historial disponibles en terreno para evitar el quiebre
        del 7,8\,\% y los bloqueos sin notificación.
  \item \textbf{Planificación algorítmica y costeo por hecho:} motor de ruteo que
        resguarda el conocimiento crítico del planificador y costeo de servir por
        cliente y por entrega construido desde los hechos operacionales.
  \item \textbf{Obtención del apoyo de los involucrados clave:} incorporación temprana
        del sindicato de conductores, del equipo de TI (4 personas) y de los operadores
        de bodega mediante validación en marcha blanca y capacitación compatible con la
        operación real.
\end{enumerate}

\subsection{Fechas e hitos relevantes}

La planificación del proyecto se estructura en torno a hitos de negocio y ventanas de
congelamiento obligatorias (Restricción R-10: sin intervención del 1 al 25 de
septiembre, todo diciembre y los primeros tres días hábiles de cada mes).

\begin{table}[htbp]
\centering
\small
\caption{Hitos de fechas relevantes}
\label{tab:hitos}
\begin{tabularx}{\textwidth}{@{}L{4.2cm}L{8.5cm}L{3cm}@{}}
\toprule
\textbf{Hito} & \textbf{Fecha / Descripción} & \textbf{Origen} \\
\midrule
PEAK de pedidos &
Septiembre: el volumen se duplica durante 3 semanas; es el período de mayor
exigencia del año y condiciona las ventanas de congelamiento &
Cap 13.2 \\
\addlinespace
Ventanas de congelamiento &
1 al 25 de septiembre, todo diciembre y primeros tres días hábiles de cada mes &
Restricción R-10 \\
\addlinespace
Desarrollo Etapa 1 &
Meses 1--12: levantamiento y línea base de alcance, diseño de arquitectura, construcción, integraciones, migración de datos, ambientes DEV, QA, PREPROD y PROD, pruebas integrales, pruebas de seguridad y certificación &
BA Art. 17.1 \\
\addlinespace
Marcha blanca Etapa 1 &
Meses 13--15: operación supervisada con datos y usuarios reales, con plan de reversión activo y medición diaria de indicadores &
BA Art. 17.1 \\
\addlinespace
Producción Etapa 1 &
Mes 16: la Etapa 1 pasa a producción y se convierte en el sistema de registro oficial del alcance comprometido &
BA Art. 17.1 \\
\addlinespace
Desarrollo Etapa 2 &
Meses 13--18: segundo alcance funcional en paralelo con la marcha blanca y la estabilización de la Etapa 1; cierre del desarrollo en el mes 18 inclusive &
BA Art. 17.1 \\
\addlinespace
Marcha blanca Etapa 2 &
Meses 19--20: operación supervisada de la Etapa 2 conviviendo con la Etapa 1 en producción, con integridad de datos entre ambos alcances &
BA Art. 17.1 \\
\addlinespace
Producción Etapa 2 &
Mes 21: la Etapa 2 pasa a producción y se acepta el PROYECTO de implementación &
BA Art. 17.1 \\
\addlinespace
Operación &
Meses 21--56: 36 meses continuos de soporte y operación de la plataforma bajo los niveles de servicio del Art. 78 &
BA Art. 17.1 \\
\bottomrule
\end{tabularx}
\end{table}

\section{Alcance de la Etapa 1 (meses 1--12)}

La Etapa 1 implementa el \textbf{primer alcance funcional y de infraestructura del PROYECTO}, de primera prioridad conforme a las Bases Técnicas del caso (BA Art.\ 15.1): comprende la \textbf{plataforma híbrida completa}, la \textbf{seguridad}, la \textbf{observabilidad}, las \textbf{integraciones críticas} y la funcionalidad que el caso declara urgente (Cap.\ 13.1). Su desarrollo ocurre entre los \textbf{meses 1 y 12}; le siguen la marcha blanca (meses 13--15) y la producción desde el \textbf{mes 16} (BA Art.\ 17.1).

\subsection{Alcance de infraestructura y plataforma (la base que reutiliza todo el PROYECTO)}

\begin{itemize}
  \item \textbf{Plataforma híbrida completa}: nube pública multi-zona (región primaria y secundaria declaradas) más componentes \emph{on-premise} en los seis sitios, definida como código (IaC), versionada y reproducible en su totalidad.
  \item \textbf{Seguridad integral}: identidad centralizada con SSO (RF-15.01--15.02), MFA para administradores y acceso externo (RF-15.03), RBAC y segregación de funciones (RF-15.04), cero confianza (\emph{Zero Trust}), matriz de controles trazable a ISO/IEC~27001 (RF-14.02, del catálogo de seguridad), plan de respuesta a incidentes y notificación de brechas (RF-14.04--14.05, del catálogo de seguridad).
  \item \textbf{Observabilidad unificada} nube + \emph{on-premise} con NOC 24$\times$7$\times$365, alertamiento por síntomas de negocio y análisis de causa raíz (RF-16.01--16.04).
  \item \textbf{Ambientes y continuidad}: habilitación de los cinco ambientes exigidos (DEV, QA, PREPROD, PROD y el ambiente de Recuperación ante Desastres, conforme a RT-04.01); recuperación ante desastres con RTO $\leq$ 4~h y RPO $\leq$ 15~min (RNF-20.06) y respaldo 3-2-1-1-0 (RNF-20.07).
  \item \textbf{Plataforma transversal de gestión}: administración de usuarios, parametrización de reglas con aprobación de dos perfiles, registro de auditoría inalterable, flujos de trabajo y bandeja de tareas, gestión documental y notificaciones multicanal (RF-17.01--17.13).
  \item \textbf{Firma electrónica avanzada} (Ley 19.799) para los actos que el caso exige (RF-18.01, RT-16.17--16.18).
  \item \textbf{Integraciones críticas}: ERP y emisión de documentos tributarios (no se reemplazan), maestro de productos y servicios de DTE, mediante capa anticorrupción (RT-05.20).
\end{itemize}

\subsection{Alcance funcional del proyecto (Etapa 1) — Lo que hará}

La funcionalidad de \textbf{primera prioridad} definida en el caso (Cap.\ 13.1) que se entrega en la Etapa 1:

\begin{itemize}
  \item \textbf{Abastecimiento y Recepción} (Épica 1, RF-01.01--01.11): validación contra OC, registro obligatorio del lote (AI~10) y del vencimiento, no conformidades con cuarentena automática, lectura GS1 y etiquetado SSCC, fichas de producto, compatibilidad térmica en \emph{putaway} e integración de la recepción con el ERP.
  \item \textbf{Gestión de Almacén e Inventario} (Épica 2, RF-02.01--02.10): topología multi-sitio configurable, consolidación de stock multi-nodo, \emph{slotting} por rotación, conteo cíclico ciego, cálculo de stock disponible para la venta, inventario en tránsito, gestión FEFO con alertas de vida útil y ficha de trazabilidad de lote. \textbf{Reemplaza el WMS de 2013} (Decisión 14, Cap.\ 16.1).
  \item \textbf{Trazabilidad y Cadena de Frío} (Épica 9, RF-09.01--09.09): trazabilidad \emph{forward} y \emph{backward} por lote bajo el estándar GS1~EPCIS, registro continuo de temperaturas de cámaras y camiones con equipo de frío, excursiones térmicas con alerta y disposición configurable, registro de lotes con control sanitario y generación del cumplimiento de cadena de frío.
  \item \textbf{Preventa Móvil} (Épica 3, RF-03.01--03.16): aplicación nativa \emph{offline-first} para los 62 preventistas (turno completo de 14~h sin señal), consulta de stock, crédito y deuda vencida en línea y sin conectividad, reserva de stock al confirmar, promociones, ventana de entrega comprometible, pedido con identificador único y descarte de duplicados al sincronizar.
  \item \textbf{Preparación de Pedidos (\emph{Picking}) y Carga} (Épica 5, RF-05.01--05.07): misiones de \emph{picking}, validación por escaneo GS1 con confirmación $\leq$ 1~s operable con guantes a $-22$\,°C, motivo obligatorio de faltante, asignación FEFO, secuenciación térmica, envío de la preparación al ERP y carga dirigida inversa a la ruta.
  \item \textbf{Transporte, Entrega y Prueba de Entrega} (Épica 6, RF-06.01--06.13): confirmación de bultos, firma digital o evidencia alternativa (QR, o fotografía con nombre del receptor sin conectividad), rechazo de productos, local cerrado con reagendamiento automático, registro de recaudación en ruta, autenticación del conductor externo por OTP, captura y saldo de envases retornables y sincronización del 100\,\% al reconectar.
  \item \textbf{Control de efectivo y cobranza en ruta} (Épica 7): rendición digital individual por camión (RF-07.02), causales tipificadas de descuadre con bloqueo del cierre (RF-07.03), alerta y bloqueo de crédito en preventa (RF-07.05), cobro con POS móvil (RF-07.06) y abonos parciales en la entrega (RF-07.10).
  \item \textbf{Logística inversa en ruta y en góndola} (Épica 8): devoluciones en la entrega (RF-08.01), destino de los productos devueltos en el andén (RF-08.02), cuenta corriente de envases retornables por cliente (RF-08.03), alertas por exceso de envases (RF-08.04) y registro de mermas por vencimiento en góndola (RF-08.05).
  \item \textbf{Núcleo del motor de ruteo} (Épica 4): asignación y secuenciación automática que codifica el conocimiento del planificador (RF-04.01), ajuste manual y recálculo de ETA (RF-04.02, RF-04.04) y parametrización de ventanas horarias por cliente (RF-04.07), para resguardar la continuidad operacional ante la jubilación del planificador (Decisión 16).
\end{itemize}

\subsection{Alcance operacional de la Etapa 1}

\begin{itemize}
  \item \textbf{Usuarios}: 62 preventistas, $\sim$200 conductores, 310 personas del centro de distribución, planificadores, calidad, administración y gerencia; integración con el ERP y la emisión de DTE.
  \item \textbf{Dispositivos}: terminales \emph{rugged} para bodega, preventa y reparto; escáneres GS1; impresoras térmicas Bluetooth; sensores de temperatura de cámaras y de camiones.
  \item \textbf{Instalaciones}: los seis sitios en cuatro regiones (centro de distribución en Talca con sala técnica secundaria de respaldo, el centro de distribución de Concepción, las tres plataformas de \emph{cross-docking} y el borde operacional del CD).
  \item \textbf{Cobertura geográfica}: centro-sur de Chile; 14.200 puntos de entrega, incluidas las rutas rurales sin cobertura que operan \emph{offline}.
  \item \textbf{Período cubierto}: desarrollo entre los meses 1--12; marcha blanca meses 13--15; producción desde el mes 16.
\end{itemize}

\section{Alcance de la Etapa 2 (meses 13--18)}

La Etapa 2 implementa el \textbf{segundo alcance funcional} definido en las Bases Técnicas del caso (BA Art.\ 15.1), \textbf{construido sobre la plataforma y los módulos de la Etapa 1 sin rehacer su arquitectura}. Su desarrollo ocurre entre los \textbf{meses 13 y 18 inclusive}, en paralelo con la marcha blanca y la estabilización de la Etapa 1 (BA Art.\ 17.2); le siguen la marcha blanca (meses 19--20) y la producción en el \textbf{mes 21}, sin degradar la disponibilidad, el desempeño ni la integridad de los datos de la Etapa 1.

\subsection{Alcance funcional del proyecto (Etapa 2) — Lo que haremos}

\begin{itemize}
  \item \textbf{Canal Moderno (EDI) y Portales} (Épica 12, RF-12.01--12.30): recepción y validación del pedido electrónico, resolución de códigos contra el maestro mediante equivalencias por cadena (RF-12.03), bandeja de excepciones (RF-12.06), confirmación o rechazo a la cadena, aviso de despacho anticipado (RF-12.09), planificación de la llegada dentro de la ventana de 30 minutos (RF-12.10), evidencia y acuse de recibo digital del canal moderno (RF-12.12--12.13), portal de clientes con registro, autenticación y consulta de entregas, documentos tributarios y estado de cuenta, portal de transportistas y portal de proveedores. La integración es configurable para nuevas cadenas (RNF-12.02) y queda en producción antes de enero de 2029 (RNF-12.01).
  \item \textbf{Analítica y Costo de Servir} (Épica 11, RF-11.01--11.08): OTIF unificado calculado a diario (RF-11.03), \emph{Fill Rate} y \emph{Perfect Order} medidos con el POD móvil (RF-11.04), costo de servir por entrega y por cliente construido desde los hechos operacionales (RF-11.02, RF-04.08), ocupación de flota (RF-11.05), tablero operacional en tiempo real (RF-11.06), liquidación y cierre comercial por camión (RF-11.07) y segmentación de clientes por rentabilidad neta (RF-11.08).
  \item \textbf{Ruteo ampliado} (Épica 4): bloqueo por exceso de capacidad (RF-04.03), restricción por requisito de cadena de frío (RF-04.05), notificación de hora estimada de llegada al cliente (RF-04.06) y cálculo integral del costo de la entrega realizada (RF-04.08).
  \item \textbf{Gestión de Flota y Telemetría} (Épica 14, RF-14.01--14.06 del catálogo funcional, sin colisión con las RF-14.* de seguridad): telemetría de camiones propios, ruta planificada vs.\ real, alerta de desviación, geocercas de llegada/salida y vinculación conductor--camión por viaje, respetando el acuerdo sindical y la privacidad del conductor (Decisión 13).
  \item \textbf{Tesorería y gestión de cartera} (Épica 7): consola de gestión de la cartera de crédito (RF-07.04), estado de cuenta en el portal de clientes (RF-07.07), validación en tesorería de los pagos web (RF-07.08), interfaz asíncrona de cobranzas hacia el ERP (RF-07.09) y cobranza presencial de facturas (RF-07.11).
  \item \textbf{Envases y mermas (analítica)} (Épica 8, RF-08.06--08.07): reporte de pérdida de envases retornables por cliente, tipo y zona, e integración de devoluciones y mermas al costo de servir.
\end{itemize}

\subsection{Justificación del reparto entre Etapas}

El reparto entre la Etapa 1 y la Etapa 2 (caso Cap.\ 17.1, \$843) se justifica por las dependencias técnicas, el riesgo, los hitos externos y la capacidad de absorción del CLIENTE:

\begin{itemize}
  \item \textbf{Dependencias técnicas}: la trazabilidad lote-a-lote exige que la recepción, el almacén/inventario, el \emph{picking} y la entrega (POD) capturen y consulten el mismo lote; la Etapa 1 los agrupa para que el ciclo completo madure junto.
  \item \textbf{Urgencia y riesgo con fecha conocida}: la trazabilidad es la habilitación comercial ante proveedores y autoridad sanitaria, y la jubilación del planificador obliga a codificar su conocimiento en el motor de ruteo durante la Etapa 1 (Decisión 16; Cap.\ 13.1--13.2).
  \item \textbf{Preferencia declarada del comité} (Cap.\ 13.1): la preventa y la entrega con evidencia digital se incorporan en la Etapa 1 para no seguir perdiendo clientes; la optimización de rutas y el costo de servir se completan en la Etapa 2.
  \item \textbf{Hitos externos} (Cap.\ 13.2): el canal moderno debe estar en producción antes de enero de 2029; con la producción de la Etapa 2 en el mes 21 queda holgado, por lo que se incorpora en la Etapa 2.
  \item \textbf{Capacidad de absorción y cronograma} (BA Art.\ 17): el segundo alcance se construye sobre la misma plataforma, lo que permite el solapamiento obligatorio de los meses 13--15 sin degradar los niveles de servicio y con una única fuente de verdad para los datos compartidos (BA Art.\ 17.2).
  \item \textbf{Criterio de asignación}: lo que depende del dato de terreno (preventa, POD, control de efectivo) entra con la Etapa 1; lo analítico y de integración externa (costo de servir, EDI, telemetría) entra con la Etapa 2, alimentándose de los datos que la Etapa 1 ya produce.
\end{itemize}

\subsection{Lo que no haremos (exclusiones del alcance)}

Lo siguiente \textbf{no formará parte} del alcance funcional de la primera ni de la segunda etapa:

\begin{itemize}
  \item El reemplazo del sistema de gestión empresarial (ERP) ni de ninguno de sus módulos,
        ni la emisión de los documentos tributarios electrónicos (los emite el ERP, con el
        que la solución se integra mediante capa anticorrupción).
  \item La administración de remuneraciones ni de recursos humanos, que residen en el
        sistema de gestión.
  \item Un sistema de punto de venta para el local del cliente, sin perjuicio del canal
        de autoatención que la solución ofrece.
  \item El comercio electrónico dirigido al consumidor final.
  \item La conectividad ni los sistemas de información de los puntos de entrega
        (clientes); la solución opera \emph{offline} cuando la red no está disponible.
  \item El desarrollo o la fabricación de hardware propio: los dispositivos se adquieren
        estándar (tipo \emph{rugged}).
  \item La obra civil y la infraestructura de climatización o de energía de los recintos
        del CLIENTE, salvo la adecuación física de la sala técnica declarada en el padrón
        de emplazamiento (Decisión 29).
\end{itemize}

\section{Enfoque para la obtención del apoyo de los involucrados clave}

La estrategia de adopción se orienta a los grupos de interés que condicionan el éxito de
la operación (detalle en capítulo de innovaciones y en la justificación con grupos
clave):

\begin{itemize}
  \item \textbf{Sindicato de conductores:} resolver la objeción a cámaras en cabina y
        control de jornada por GPS (Restricción R-11), proponiendo mecanismos de
        geolocalización exclusiva para control operativo y logístico (tiempos de
        arribo), no para supervisión de jornada.
  \item \textbf{Equipo de TI del cliente (4 personas):} todo lo que requiera un
        administrador dedicado se ofrece como servicio y costeado (Restricción R-08).
  \item \textbf{Bodega (turno de noche, rotación 38\,\%):} soluciones que no requieran
        entrenamiento prolongado (Restricción R-09), hardware \emph{rugged} operable con
        guantes a $-$22\,°C y validación en marcha blanca.
  \item \textbf{Transportistas externos:} mecanismos ágiles de incorporación de
        conductores que rotan sin aviso (Restricción R-06).
\end{itemize}

\section{Reglas de negocio}

Conforme al Capítulo 17.1 del caso, el proyecto mantiene un \textbf{registro de reglas de
negocio}: las reglas de la industria que la solución debe respetar y que el caso no
explícita como requerimiento formal. Cada regla indica qué se captura, en qué punto del
proceso, por quién, con qué consecuencia si se incumple, el origen en el caso y el
requerimiento RF/RNF que la implementa. Las reglas se derivan de las decisiones del
numeral 16.1 ya resueltas (véase el registro de supuestos), de modo que no contradicen la
arquitectura ni el alcance propuestos.

\begin{table}[htbp]
\centering
\small
\caption{Trazabilidad de reglas de negocio}
\label{tab:reglas_resumen}
\begin{tabularx}{\textwidth}{@{}L{1.3cm}L{4.1cm}L{1.3cm}L{7.2cm}@{}}
\toprule
\textbf{Código} & \textbf{Regla} & \textbf{Decisión} & \textbf{RF/RNF principal} \\
\midrule
RNG-01 & Reserva de stock (orden de confirmación) & D8 & RF-03.03, RF-02.05 \\
RNG-02 & Stock disponible vs.\ comprometido / en tránsito & D8 & RF-02.05, RF-02.06 \\
RNG-03 & Política de crédito en preventa & D8/D40 & RF-03.04--03.10 \\
RNG-04 & Excursión térmica y disposición & D4 & RF-09.05--09.07, RF-01.11 \\
RNG-05 & Local cerrado / reintento & D3 & RF-06.04, RF-04.07 \\
RNG-06 & Devoluciones y efecto sobre documento tributario & D12 & RF-06.03, RF-08.01/08.02 \\
RNG-07 & Control de envases retornables & D10 & RF-08.03/08.04/08.06 \\
RNG-08 & Cambio de precio (toma vs.\ despacho) & D9 & RF-03.11--03.13 \\
RNG-09 & Entrega cumplida / OTIF y Fill Rate & D1 & RF-11.03, RF-11.04 \\
RNG-10 & FEFO y vida útil remanente & --- & RF-05.05, RF-02.09 \\
RNG-11 & Maestro de productos y códigos GS1 & D11 & RF-01.10, RF-12.03/12.06 \\
RNG-12 & Sincronización y resolución de conflictos & RT-03.12 & RF-02.05, RF-03.16/17 \\
RNG-13 & Secuenciación térmica / validación por zona & --- & RF-05.04, RF-01.11/01.08 \\
RNG-14 & Capacidad del vehículo y carga inversa & --- & RF-04.03/04.05, RF-05.07 \\
\bottomrule
\end{tabularx}
\vspace{4pt}
{\small\color{lafroxGrato}\pendiente{} Completar el detalle de cada regla en el
registro de requerimientos del Cap.\ 17.1.}
\end{table}

\subsubsection{Reglas de inventario y stock}

\begin{table}[htbp]
\centering
\small
\begin{tabularx}{\textwidth}{@{}L{2.4cm}L{8.3cm}L{4.2cm}@{}}
\toprule
\textbf{RNG-01} \textbar{} Asignación y reserva de stock &
\textbf{Qué:} el momento en que el stock se compromete y el orden ante pedidos que
compten por el mismo producto. \textbf{Punto:} al confirmar el pedido en el sistema de
gestión. \textbf{Por quién:} sistema en la confirmación; el preventista gatilla; el jefe
de bodega administra. \textbf{Regla:} reserva por orden de llegada; el segundo
confirmado queda con quiebre parcial registrado. \textbf{Consecuencia:} doble compromiso
de stock, quiebres de despacho y OTIF caído. &
\textbf{Origen:} Cap.\ 16.1 D8. \textbf{RF:} RF-03.03, RF-02.05, RF-05.03. \\
\addlinespace
\textbf{RNG-02} \textbar{} Stock disponible vs.\ comprometido y en tránsito &
\textbf{Qué:} definición de stock disponible y tratamiento del inventario en tránsito.
\textbf{Punto:} consulta de stock (preventa) y consolidación multi-sitio (bodega).
\textbf{Por quién:} sistema; preventista y jefe de bodega. \textbf{Regla:} disponible =
físico -- comprometido; el stock en tránsito no suma hasta la recepción confirmada; en
offline el dato es indicativo con antigüedad, y los conflictos se resuelven por orden
cronológico. \textbf{Consecuencia:} sobreventa o venta de inventario en tránsito. &
\textbf{Origen:} Cap.\ 2.3, D8; RT-03.12. \textbf{RF:} RF-02.05, RF-02.06, RF-03.02. \\
\addlinespace
\textbf{RNG-10} \textbar{} FEFO y vida útil remanente &
\textbf{Qué:} orden de salida del inventario y alerta por vencimiento. \textbf{Punto:}
asignación de picking y alarma de vida útil. \textbf{Por quién:} sistema dirige; el
preparador confirma; el jefe de bodega gestiona alertas. \textbf{Regla:} salida
estrictamente FEFO con bitácora de excepción; alerta al superar el umbral de vida útil
remanente (p.\,ej.\ 30 días). \textbf{Consecuencia:} vencimiento en almacén y merma. &
\textbf{Origen:} Cap.\ 4.8, 9.10. \textbf{RF:} RF-05.05, RF-02.09, RF-05.03. \\
\bottomrule
\end{tabularx}
\end{table}

\subsubsection{Reglas de crédito}

\begin{table}[htbp]
\centering
\small
\begin{tabularx}{\textwidth}{@{}L{2.4cm}L{8.3cm}L{4.2cm}@{}}
\toprule
\textbf{RNG-03} \textbar{} Política de crédito en preventa &
\textbf{Qué:} límite de crédito del cliente y comportamiento ante su superación.
\textbf{Punto:} toma de pedido (en línea y offline). \textbf{Por quién:} el preventista
consulta; el supervisor de crédito autoriza excepciones; el sistema aplica bloqueos.
\textbf{Regla:} superar el límite dispara alerta; superar el 110\,\% bloquea el envío;
toda excepción se autoriza de forma registrada; en offline la excepción queda
aprobada a la sincronización con respaldo presencial. \textbf{Consecuencia:} ventas a
clientes morosos y descuadre. &
\textbf{Origen:} Cap.\ 16.1 D8/D40. \textbf{RF:} RF-03.04/03.05, RF-03.08,
RF-03.09, RF-03.10. \\
\bottomrule
\end{tabularx}
\end{table}

\subsubsection{Reglas de calidad y cadena de frío}

\begin{table}[htbp]
\centering
\small
\begin{tabularx}{\textwidth}{@{}L{2.4cm}L{8.3cm}L{4.2cm}@{}}
\toprule
\textbf{RNG-04} \textbar{} Excursión térmica y disposición &
\textbf{Qué:} qué constituye una excursión y quién decide la disposición. \textbf{Punto:}
transporte y cámaras; bloqueo de despacho. \textbf{Por quién:} jefa de calidad define
parámetros; el sistema alerta; el conductor decide la disposición final. \textbf{Regla:}
excursión por parámetros configurables (tiempo y grados); alerta en tiempo real; el
bloqueo del despacho NO es automático (el sistema habilita y el conductor decide);
monitorear $\neq$ decidir (HACCP/GDP). \textbf{Consecuencia:} riesgo sanitario y retiro. &
\textbf{Origen:} Cap.\ 16.1 D4, Cap.\ 10; RT-17.06. \textbf{RF:} RF-09.05, RF-09.06,
RF-09.07, RF-01.11. \\
\addlinespace
\textbf{RNG-13} \textbar{} Secuenciación térmica y validación por zona &
\textbf{Qué:} orden de recolección y validación de zona térmica. \textbf{Punto:} picking
y putaway. \textbf{Por quién:} sistema secuencia; preparador confirma; operador valida.
\textbf{Regla:} misión secuenciada seco $\rightarrow$ refrigerado $\rightarrow$ congelado;
se recoge el 100\,\% de ambiente seco antes de liberar cámaras; validation por
compatibilidad térmica en putaway (un congelado no se ubica fuera de zona congelada).
\textbf{Consecuencia:} ruptura de cadena de frío y rechazo. &
\textbf{Origen:} Cap.\ 4.5, 2.3. \textbf{RF:} RF-05.04, RF-01.11, RF-01.08. \\
\bottomrule
\end{tabularx}
\end{table}

\subsubsection{Reglas de reparto y entrega}

\begin{table}[htbp]
\centering
\small
\begin{tabularx}{\textwidth}{@{}L{2.4cm}L{8.3cm}L{4.2cm}@{}}
\toprule
\textbf{RNG-05} \textbar{} Reintento de entrega / local cerrado &
\textbf{Qué:} conducta ante un local cerrado en la ventana comprometida. \textbf{Punto:}
la entrega. \textbf{Por quién:} el conductor reporta; el sistema reagenda
automáticamente; el planificador administra. \textbf{Regla:} si el local está cerrado, se
registra el envío fallido y se reagenda a la ventana siguiente registrada; el conductor
no decide por su cuenta; si persiste, los productos regresan al almacén.
\textbf{Consecuencia:} reentregas y OTIF bajo. &
\textbf{Origen:} Cap.\ 16.1 D3. \textbf{RF:} RF-06.04, RF-04.07, RF-04.06. \\
\addlinespace
\textbf{RNG-09} \textbar{} Entrega cumplida / OTIF y Fill Rate &
\textbf{Qué:} criterio de entrega cumplida y métricas de servicio. \textbf{Punto:}
medición diaria del servicio. \textbf{Por quién:} gerente comercial; sistema registra el
POD; jefa de calidad. \textbf{Regla:} la entrega parcial NO cuenta como cumplida (In-Full
= 100\,\%); On-Time = fecha y ventana pactada; Fill Rate = unidades entregadas vs.\
pedidas; Perfect Order contra la orden original; la parcial queda como desvío auditado.
\textbf{Consecuencia:} imposibilidad de acordar el estado real del servicio. &
\textbf{Origen:} Cap.\ 16.1 D1; Cap.\ 14.2. \textbf{RF:} RF-11.03, RF-11.04. \\
\addlinespace
\textbf{RNG-14} \textbar{} Capacidad del vehículo y carga inversa &
\textbf{Qué:} límites de capacidad y orden de carga. \textbf{Punto:} planificación de
rutas y carga del camión. \textbf{Por quién:} planificador; sistema bloquea; cargador
ejecuta. \textbf{Regla:} no se asigna línea que exceda peso/volumen nominal; no se admite
cadena de frío en vehículo seco; la carga se dirige en secuencia inversa a la ruta y al
completar se verifica que todos los bultos fueron escaneados. \textbf{Consecuencia:}
sobrecarga, rechazo y pérdida de tiempo de descarga. &
\textbf{Origen:} Cap.\ 4.4--4.5. \textbf{RF:} RF-04.03, RF-04.05, RF-05.07. \\
\bottomrule
\end{tabularx}
\end{table}

\subsubsection{Reglas de devoluciones, envases y maestros}

\begin{table}[htbp]
\centering
\small
\begin{tabularx}{\textwidth}{@{}L{2.4cm}L{8.3cm}L{4.2cm}@{}}
\toprule
\textbf{RNG-06} \textbar{} Devoluciones y efecto sobre documento tributario &
\textbf{Qué:} la devolución en la entrega y su efecto tributario. \textbf{Punto:} POD y
destino en el andén de retorno. \textbf{Por quién:} conductor registra; jefe de bodega
decide el destino; finanzas ejecuta la nota de crédito. \textbf{Regla:} se registra en el
POD (total/parcial, con SKU, cantidad y motivo), operando offline; el destino se decide
en el andén y se transmite al ERP; el efecto sobre el DTE se resuelve con nota de crédito
del ERP (que no se reemplaza), con retención legal. \textbf{Consecuencia:} descuadre
contable y pérdida de trazabilidad. &
\textbf{Origen:} Cap.\ 16.1 D12; RT-16.14. \textbf{RF:} RF-06.03, RF-08.01, RF-08.02;
RNF-08.02. \\
\addlinespace
\textbf{RNG-07} \textbar{} Control de envases retornables &
\textbf{Qué:} mecanismo de control del parque retornable. \textbf{Punto:} cada entrega
(descuento) y devolución (suma). \textbf{Por quién:} conductor captura; el sistema lleva
la cuenta corriente por cliente; el preventista recibe alertas. \textbf{Regla:} control
por saldo por cliente (no por unidad identificada); se alerta al superar el umbral y se
reporta la pérdida por cliente, tipo y zona. \textbf{Consecuencia:} pérdida no controlada
del parque. &
\textbf{Origen:} Cap.\ 16.1 D10; Cap.\ 14.1. \textbf{RF:} RF-08.03, RF-08.04, RF-08.06,
RF-06.09. \\
\addlinespace
\textbf{RNG-08} \textbar{} Cambio de precio (toma vs.\ despacho) &
\textbf{Qué:} precio aplicable cuando cambia entre toma y despacho. \textbf{Punto:}
facturación / despacho de la línea. \textbf{Por quién:} administrador comercial configura
la excepción; el sistema aplica el por defecto; el cliente recibe la alerta.
\textbf{Regla:} por defecto se cobra el precio vigente al despacho; solo con excepción
explícita se respeta el precio de la toma; sin excepción y con cambio, se alerta al
cliente antes de o con la facturación. \textbf{Consecuencia:} reclamos y litigios de
facturación. &
\textbf{Origen:} Cap.\ 16.1 D9. \textbf{RF:} RF-03.11, RF-03.12, RF-03.13. \\
\addlinespace
\textbf{RNG-11} \textbar{} Maestro de productos y códigos GS1 &
\textbf{Qué:} qué maestro manda ante cambios de código y cómo se resuelven los códigos de
cadenas. \textbf{Punto:} ficha de producto y recepción EDI. \textbf{Por quién:}
administrador de catálogo; sistema resuelve equivalencias; operador de excepciones EDI.
\textbf{Regla:} prevalece el maestro interno; la ficha es la única fuente autorizada con
historial de cambios; todo código externo se resuelve por tabla de equivalencias por
cadena; lo no resuelto queda en bandeja de excepciones. \textbf{Consecuencia:} producto
duplicado o pedidos EDI rechazados. &
\textbf{Origen:} Cap.\ 16.1 D11; Cap.\ 16.2 (GS1). \textbf{RF:} RF-01.10, RF-12.03,
RF-12.06. \\
\bottomrule
\end{tabularx}
\end{table}

\subsubsection{Reglas de integración y sincronización}

\begin{table}[htbp]
\centering
\small
\begin{tabularx}{\textwidth}{@{}L{2.4cm}L{8.3cm}L{4.2cm}@{}}
\toprule
\textbf{RNG-12} \textbar{} Enlace, sincronización y resolución de conflictos &
\textbf{Qué:} comportamiento tras la reconexión y el criterio de resolución de
conflictos. \textbf{Punto:} toda operación desconectada (CD 24 h, terreno 14 h) y su
reconciliación. \textbf{Por quién:} sistema de sincronización automática; bitácora
auditable. \textbf{Regla:} al restablecer el enlace, sincronización automática con
reconciliación determinista y bitácora auditable; en terreno (14 h) el dato offline es
indicativo con antigüedad y se valida contra el servidor al reconectar.
\textbf{Consecuencia:} pedidos perdidos o duplicados y quiebres no reconciliados. &
\textbf{Origen:} RT-03.12; Cap.\ 10. \textbf{RF:} RF-02.05, RF-03.16/17; RNF-13.01/13.02. \\
\bottomrule
\end{tabularx}
\end{table}

\section{Enlace al catálogo de requerimientos}

El alcance funcional se despliega en el catálogo de requerimientos funcionales (RF),
no funcionales (RNF) y de bases, trazable a su origen en el caso:

\begin{itemize}
  \item Requerimientos funcionales: \emph{Requerimientos LOGISTICA -- Consolidado de
        Requerimientos Funcionales}.
  \item Requerimientos no funcionales: \emph{Requerimientos LOGISTICA -- Requerimientos
        NO Funcionales}.
  \item Requerimientos de bases: \emph{Requerimientos LOGISTICA -- Requerimientos de
        bases}.
\item Planilla consolidada:
        \url{https://docs.google.com/spreadsheets/d/19HYOFHX_maoNcECLjazVm9j3hZP-iplyxgyHAy7HDuo}
\end{itemize}

\section{Cat\'alogos de requerimientos}

Esta secci\'on consolida los cat\'alogos completos de requerimientos funcionales
(RF), no funcionales (RNF) y de bases (RT) del proyecto, seg\'un el registro de
requerimientos del cap\'itulo 17.1 del caso. Cada requerimiento lista su
identificador, t\'itulo, actor, descripci\'on, prioridad y origen.

\subsection{Requerimientos funcionales (RF)}
\begin{longtable}{p{1.7cm} p{2.3cm} p{1.9cm} p{5.0cm} p{1.1cm} p{1.9cm}}
\caption{Cat\'alogo de requerimientos funcionales (RF)}\label{tab:rf}\\
\toprule
\textbf{ID} & \textbf{T\'itulo} & \textbf{Actor} & \textbf{Descripci\'on} & \textbf{Prioridad} & \textbf{Origen} \\
\midrule
\endfirsthead
\toprule
\textbf{ID} & \textbf{T\'itulo} & \textbf{Actor} & \textbf{Descripci\'on} & \textbf{Prioridad} & \textbf{Origen} \\
\midrule
\endhead
\bottomrule
\endlastfoot
\textbf{RF-01.01} & Recepción y Validación de Cantidades contra OC & Recepcionista de bodega & El sistema debe registrar la recepción física de un proveedor validando las cantidades recibidas por SKU contra la OC y generando una alerta por cada discrepancia, sin bloquear el cierre. & Muy Alta & Cap. 4.1, Anexo A, Cap 12 \\
\textbf{RF-01.02} & Registro Obligatorio de Lote & Recepcionista de bodega & El sistema debe exigir el número de lote del proveedor en campo estructurado para todo SKU con trazabilidad obligatoria, capturándolo por escaneo GS1 (AI 10) o ingreso manual, y bloqueando el avance si el campo está vacío. & Muy Alta & Cap. 4.1, Cap 4.9, Cap. 7.3, \\
\textbf{RF-01.03} & Registro de Fecha de Vencimiento & Recepcionista de bodega & El sistema debe registrar la fecha de vencimiento en formato DD/MM/AAAA para todo SKU con control de vencimiento habilitado, vinculándola al lote y a la ubicación asignada como entrada para la lógica FEFO. & Muy Alta & Cap. 4.1, Cap. 4.8 \\
\textbf{RF-01.04} & Registro de No Conformidades en Recepción & Recepcionista de bodega & El sistema debe permitir registrar una no conformidad durante la recepción seleccionando su tipo desde una lista predefinida e ingresando SKU, lote, cantidad afectada y al menos una fotografía capturada desde el dispositivo. & Alta & Cap. 4.1, Cap. 4.9, Cap 8 (Entrevista de Calidad), RT-17.06 \\
\textbf{RF-01.05} & Cuarentena Automática por No Conformidad & Recepcionista de bodega & Al confirmar una no conformidad, el sistema debe cambiar automáticamente a estado ‘Cuarentena’ la cantidad afectada del SKU/lote correspondiente, bloqueándola para picking, y notificar al perfil de calidad mediante una alerta en la aplicación. & Alta & Cap. 4.1, Cap. 4.9, RT-17.06 \\
\textbf{RF-01.06} & Lectura y Decodificación de Códigos GS1 & Recepcionista de bodega & El sistema debe decodificar códigos de barras GS1 interpretando los AIs 01 (GTIN), 10 (lote), 17 (vencimiento) y 00 (SSCC), autocompletando los campos en pantalla; ante código inválido o no registrado, habilita el ingreso manual. & Alta & RT-17.06, Cap. 16.2 (GS1) \\
\textbf{RF-01.07} & Generación e Impresión de Etiqueta SSCC & Recepcionista de bodega & Al finalizar el registro de un pallet, el sistema debe generar automáticamente un SSCC bajo estándar GS1-128 vinculado al GTIN, lote, fecha de vencimiento y fecha de recepción, e imprimir la etiqueta en la impresora térmica del andén. & Media & Cap. 4.2 , Cap 12 (GS1), RT-05.23 \\
\textbf{RF-01.08} & Asignación de SSCC a Ubicación de Almacenamiento & Operador de bodega & El sistema debe permitir asociar un SSCC a una dirección de almacenamiento (pasillo-columna-nivel) escaneando ambas etiquetas y validando la compatibilidad térmica antes de confirmar. & Media & Cap. 4.2, Cap. 12 (GS1), RT-05.23 \\
\textbf{RF-01.09} & Integración de Recepción con ERP & Recepcionista de bodega & Al cerrar una recepción, el sistema debe enviar automáticamente la transacción de ingreso de inventario al ERP mediante su interfaz de integración; si el ERP no está disponible, encola la transacción para reintento automático. & Muy Alta & Cap. 5, Anexo A, RT-03.13 \\
\textbf{RF-01.10} & Gestión de Fichas de Producto & Administrador de catálogo / Jefe de abastecimiento & El sistema debe permitir crear y mantener fichas de producto con atributos obligatorios: descripción, formato, GTIN, SKU interno, tipo de almacenamiento, temperatura mínima y máxima (\ensuremath{^\circ}C), peso neto (kg), dimensiones (cm), días de vida útil, indicadores de trazabilidad y control de vencimiento (sí/no), clasificación sanitaria y RUT del proveedor. & Muy Alta & Cap 14.1 (pasa de 8.400 a 9.500 SKu), Cap 8 (Hugo), Cap 16.1 (N11) \\
\textbf{RF-01.11} & Validación de Compatibilidad Térmica en Putaway & Operador de bodega & El sistema debe validar la compatibilidad térmica entre el tipo de almacenamiento del SKU y la zona de temperatura de la ubicación destino al realizar un putaway, bloqueando la operación si no son compatibles. & Muy Alta & Cap. 2.3, Cap. 4.9, Cap. 10 (restricción 2) \\
\textbf{RF-02.01} & Configuración de topología de instalaciones & Administrador del sistema / Jefe de bodega & El sistema debe permitir registrar y configurar cada instalación de la red logística (centros de distribución y plataformas de cross-docking), definiendo su topología interna: zonas operativas (seco con rango 15–25 \ensuremath{^\circ}C, refrigerado con rango 2–8 \ensuremath{^\circ}C, congelado con rango $-$18 a $-$25 \ensuremath{^\circ}C), pasillos, racks, niveles y posiciones individuales de almacenamiento. La configuración debe estar disponible en modo offline. & Muy Alta & Cap. 2.3, RT-02.12 \\
\textbf{RF-02.02} & Consolidación de stock físico multi-sitio & Jefe de bodega / Gerente de operaciones & El sistema debe consolidar el stock físico en unidades de cada SKU en todos los nodos de la red logística. & Alta & Cap. 2.3 \\
\textbf{RF-02.03} & Reasignación de ubicaciones (slotting) & Jefe de bodega & El sistema debe calcular periódicamente una propuesta de reasignación de ubicaciones (slotting) basada en la clasificación de rotación (ABC) del último período configurable, considerando además peso, volumen y compatibilidad de zona. Los productos de clasificación A deben proponerse en ubicaciones cercanas a los andenes de salida. La propuesta debe ser aprobada por el jefe de bodega antes de generar instrucciones de movimiento. & Muy Alta & Cap. 8 (Ximena) \\
\textbf{RF-02.04} & Conteo cíclico en modo ciego & Operario de bodega / Jefe de bodega & El sistema debe generar un plan de conteo cíclico mensual que cubra una cantidad de posiciones definidas, seleccionando posiciones por criterios configurables (rotación, valor, discrepancia histórica). Las tareas de conteo deben asignarse a los dispositivos de los operarios de bodega en modo de conteo ciego, sin mostrar la cantidad teórica. Tras el ingreso de la cantidad física, el sistema compara contra el teórico y registra la diferencia con su causa en campo estructurado. & Alta & Cap. 4.2, Cap. 14.1 \\
\textbf{RF-02.05} & Cálculo de stock disponible para venta & Preventista / Jefe de bodega & El sistema debe calcular el stock disponible para venta por SKU y por centro de distribución, definido como la diferencia entre el stock físico registrado y el inventario comprometido por pedidos confirmados pendientes de preparación. En modo desconectado, la consulta debe reflejar el stock al momento de la última sincronización, indicando la antigüedad del dato. & Alta & Cap. 16.1 (Decisión \#8) \\
\textbf{RF-02.06} & Gestión de inventario en tránsito & Jefe de bodega / Operaciones & El sistema debe registrar como «en tránsito» todo inventario despachado desde un nodo de origen hacia un nodo destino, desde la carga confirmada hasta la recepción confirmada en destino, y generar la transacción de transferencia correspondiente en el ERP. & Media & Cap. 2.3 (abastecimiento desde Talca) \\
\textbf{RF-02.07} & Consulta de nivel de ocupación por zona & Jefe de bodega / Gerente de operaciones & El sistema debe permitir consultar el nivel de ocupación de cada centro de distribución, desglosado por zona operativa (seco, refrigerado, congelado), expresado como porcentaje de ubicaciones físicas utilizadas respecto de la capacidad total configurada. &  &  \\
\textbf{RF-02.08} & Gestión FEFO y alertas de vida útil & Jefe de bodega & El sistema debe dirigir la preparación de pedidos (picking) asignando tareas de recolección al dispositivo del preparador, indicando: ubicación, producto, lote, cantidad, y secuencia optimizada de recorrido. El preparador confirma cada línea en el dispositivo. Al no encontrar un producto o encontrar cantidad insuficiente, el preparador debe registrar el faltante con su causa en un campo estructurado (sin stock, producto dañado, producto vencido, error de ubicación, otro). &  &  \\
\textbf{RF-02.09} &  & Jefe de bodega & El sistema debe gestionar el almacenamiento de productos con fecha de vencimiento aplicando la regla FEFO (First Expired, First Out), priorizando en la preparación de pedidos los lotes con fecha de vencimiento más próxima. El sistema debe alertar cuando un producto en stock alcance un umbral de vida útil remanente configurable. &  &  \\
\textbf{RF-02.10} & Ficha de trazabilidad de lote & Perfil de calidad / Jefe de bodega & El sistema debe proveer una interfaz de trazabilidad que permita buscar y detallar el ciclo de vida completo de un lote. Al consultar un lote específico, debe desglosar su estado actual, ubicación física en bodega, datos de recepción de proveedor y detalle de salida (clientes/mermas). & Muy Alta & Cap. 1, Cap. 4.9, Cap. 9.5 \\
\textbf{RF-03.01} & Consulta de stock disponible en línea & Preventista & La app debe mostrar, con conectividad, el stock disponible en tiempo real para cada SKU al tomar el pedido. & Muy Alta &  \\
\textbf{RF-03.02} & Consulta de stock disponible sin conectividad & Preventista & La app debe mostrar, sin conectividad, el stock según la última sincronización válida, con su antigüedad visible. & Muy Alta &  \\
\textbf{RF-03.03} & Reserva de stock al confirmar el pedido & Preventista & El sistema debe reservar el stock comprometido en el momento de la confirmación del pedido en el sistema de gestión, por orden de llegada. & Alta &  \\
\textbf{RF-03.04} & Consulta de crédito y deuda vencida en línea & Preventista & La app debe mostrar, con conectividad, el crédito disponible y la deuda vencida del cliente al tomar el pedido. & Muy Alta &  \\
\textbf{RF-03.05} & Consulta de crédito y deuda vencida sin conectividad & Preventista & La app debe mostrar, sin conectividad, el crédito y la deuda según la última sincronización válida, con su antigüedad visible. & Muy Alta &  \\
\textbf{RF-03.06} & Consulta de historial de compra del cliente & Preventista & La app debe mostrar el historial de compra del cliente (productos y frecuencia). & Alta &  \\
\textbf{RF-03.07} & Aplicación de promociones a una línea de pedido & Preventista & La app debe permitir aplicar una promoción vigente (3x2, descuento por volumen) a una línea de pedido, calculando el descuento automáticamente. & Alta &  \\
\textbf{RF-03.08} & Alerta por superación del límite de crédito & Preventista & El sistema debe alertar al preventista cuando el pedido, sumado a la deuda vencida, supere el límite de crédito del cliente. & Muy Alta &  \\
\textbf{RF-03.09} & Bloqueo de envío por exceso de crédito & Preventista & El sistema debe bloquear el envío del pedido cuando este, sumado a la deuda vencida, supere el límite de crédito en más de [110] \%. & Muy Alta &  \\
\textbf{RF-03.10} & Autorización de excepción por supervisor de crédito & Supervisor de credito & El sistema debe permitir que un supervisor autorice, de forma registrada, el envío de un pedido bloqueado. & Alta &  \\
\textbf{RF-03.11} & Aplicación del precio vigente al despacho por defecto & Administrador comercial & El sistema debe aplicar, por defecto, el precio vigente al despacho cuando el precio cambió desde la toma del pedido. & Alta &  \\
\textbf{RF-03.12} & Configuración de excepción de precio por cliente o cadena & Administrador comercial & El sistema debe permitir configurar, por cliente o cadena, que se respete el precio de toma de pedido en lugar del precio de despacho. &  &  \\
\textbf{RF-03.13} & Alerta al cliente por cambio de precio & Cliente & El sistema debe generar una alerta al cliente cuando el precio de despacho difiera del precio pactado, si no tiene excepción configurada (RF-03.12). &  &  \\
\textbf{RF-03.14} & Información de ventana de entrega comprometible &  &  &  &  \\
\textbf{considerando stock} & capacidad de despacho y restricciones del cliente.,El módulo de planificación de rutas expone capacidad disponible (dependencia externa).,El preventista visualiza la ventana disponible antes de confirmar; la promesa queda registrada.,Muy Alta, &  &  &  &  \\
\textbf{RF-03.15} & Sugerencia de productos al preventista & Preventista & El sistema debe generar sugerencias de productos usando historial de compra, frecuencia, productos habituales y promociones vigentes. & Alta &  \\
\textbf{RF-03.16} & Generación de identificador único de pedido offline & Aplicacion de preventa & La app debe generar un identificador único en el dispositivo al crear un pedido sin conectividad. & Muy Alta &  \\
\textbf{RF-03.17} & Detección y descarte de pedidos duplicados al sincronizar & Preventista & El sistema debe detectar y descartar automáticamente los reenvíos duplicados de un pedido con el mismo identificador único al sincronizar. & Muy Alta &  \\
\textbf{RF-04.01} & Asignación y secuenciación automática de rutas & Planificador & El sistema debe asignar pedidos a la flota, secuenciando la ruta en función de la capacidad volumétrica y zonas logísticas de entrega. & Alta & Cap 4.4, Criterio de aceptación \#7 \\
\textbf{RF-04.02} & Modificación manual de rutas y reasignación de pedidos & Planificador & El sistema debe permitir modificar manualmente la secuencia de ruta y reasignar pedidos entre vehículos & Media &  \\
\textbf{RF-04.03} & Bloqueo de asignación por exceso de capacidad & Planificador / Sistema & El sistema debe bloquear la adición de líneas de pedido a una ruta si la sumatoria excede los límites de peso o volumen nominales del vehículo. & Alta &  \\
\textbf{RF-04.04} & Ajuste de secuencia de ruta por ventanas horarias & Sistema & El sistema debe permitir la configuración de ventanas horarias específicas por cliente y ajustar la secuencia para cumplirlas. & Alta &  \\
\textbf{RF-04.05} & Restricción de vehículos por requisito de cadena de frío & Planificador / Sistema & El sistema debe restringir la asignación de pedidos refrigerados o congelados exclusivamente a vehículos con equipo de frío. & Alta &  \\
\textbf{RF-04.06} & Notificación de hora estimada de llegada al cliente & Sistema & El sistema debe crear un aviso al contacto del cliente con la hora estimada de llegada incluyendo una ventana de 30 minutos. & Media &  \\
\textbf{RF-04.07} & Parametrización de ventanas horarias & Planificador / Sistema & El sistema debe permitir parametrizar franjas horarias de recepción en la ficha maestra de cada cliente (Ventana Horaria). & Alta &  \\
\textbf{RF-04.08} & Cálculo integral del costo de entrega realizada & Sistema & El sistema debe calcular el costo de cada entrega tras haber sido realizada. & Alta &  \\
\textbf{RF-05.01} & Generación de misiones de picking & Jefe de Bodega & El jefe de bodega debe poder iniciar la generación manualmente o configurar la generación automática al cierre de pedidos del turno. Las misiones deben permitir agrupación por pedido individual o por consolidación de pedidos de una misma ruta. & Muy Alta & Cap. 4.5 \\
\textbf{RF-05.02} & Validación de líneas por escaneo GS1 & Preparador & El sistema debe exigir la validación de cada línea recolectada escaneando el código GS1. El escaneo debe validar el SKU y capturar el Lote exacto tomado. Permite ingreso manual solo como contingencia auditable. & Muy Alta & Cap. 4.5, Cap. 12 \\
\textbf{RF-05.03} & Motivo obligatorio de faltante & Preparador & El sistema debe exigir la selección obligatoria de un motivo parametrizable de una lista estructurada, cada vez que se reporte una cantidad recolectada inferior a la solicitada. & Muy Alta & Cap. 8, Crit. \#15 \\
\textbf{RF-05.04} & Secuenciación térmica de misiones & Sistema & El sistema debe secuenciar las líneas de la misión de preparación priorizando las zonas en el siguiente orden: seco $\rightarrow$ refrigerado $\rightarrow$ congelado, para minimizar el tiempo de exposición de los productos de frío fuera de su zona térmica. & Alta & Cap. 4.5 \\
\textbf{RF-05.05} & Asignación de inventario FEFO en picking & Preparador & El sistema debe asignar el inventario para la preparación de pedidos aplicando estrictamente la regla lógica FEFO (First Expired, First Out) según la fecha de vencimiento del lote. & Muy Alta & Cap. 4.8, 9.10 \\
\textbf{RF-05.06} & Envío de preparación al ERP &  & El sistema debe enviar al ERP, al cerrar una misión de preparación, el detalle consolidado por cada línea: SKU, cantidad efectivamente recolectada y lote correspondiente, vinculado al ID del pedido original. & Muy Alta & Cap. 4.5 (segundo punto de pérdida) \\
\textbf{RF-05.07} & Carga dirigida inversa a ruta & Cargador / Sistema & El sistema debe dirigir la carga de pedidos preparados en el camión siguiendo la secuencia inversa de la ruta asignada, garantizando que el primer cliente a visitar sea el último cargado (más accesible). El cargador debe confirmar cada bulto/pallet escaneando su SSCC. & Alta & Cap. 4.5, Criterio de Aceptación \#14 \\
\textbf{RF-06.01} & Confirmación de recepción de bultos & Conductor & El sistema debe permitir marcar la recepción conforme de bultos mediante validación por línea de pedido individual o marcado masivo por guía. & Alta &  \\
\textbf{RF-06.02} & Captura de firma digital & Conductor / Cliente & El sistema debe permitir capturar la firma digital del receptor como evidencia de entrega, asociada a la guía de despacho electrónica. & Muy Alta &  \\
\textbf{RF-06.03} & Reporte de rechazo de productos & Conductor & El sistema debe permitir reportar un rechazo total o parcial seleccionando la cantidad y el motivo desde un catálogo predefinido. & Media &  \\
\textbf{RF-06.04} & Reporte de local cerrado y reagendamiento automático & Conductor & El sistema debe permitir reportar un local cerrado, indicando la dirección del cliente. & Media &  \\
\textbf{RF-06.05} & Registro de recaudación de pagos en efectivo & Conductor & El sistema debe permitir digitar el monto exacto recaudado en pagos realizados en efectivo & Muy Alta &  \\
\textbf{RF-06.07} & Sincronización automática de registros de turno & Conductor repartidor, Jefe de bodega / Gerente de operaciones & El sistema debe sincronizarse automaticamente, al detectar la reconexión de señal o el regreso al centro de distribución, sincronizar automaticamente todos los registros generados durante el turno, recuperando la totalidad de los despachos registrados, consolidandolos en un reporte de cierre. & Muy Alta & Cap. 4.6 y 4.7 \\
\textbf{RF-06.08} & Autenticacion conductor externo mediante OTP & Conductor repartidor y Jefe de flota / Transportista & El sistema debe permitir al conductor de una empresa externa iniciar su jornada usando las credenciales temporales por OPT definido en el RF-15.02. sin necesidad de correo corporativo. & Alta & Restricción \#6, RT-12.11, Cap. 16.1 (Decisión \#5) \\
\textbf{RF-06.09} & Captura de devolución de envaces retornables & Conductor repartidor & Al registrar una entrega, el sistema debe permitir al conductor capturar la cantidad de canastillos y pallets devueltos por el cliente, identificando el tipo y estado. & Alta & Cap. 4.8, Cap. 9.7, Criterio Aceptación \#12 \\
\textbf{RF-06.10} & Actualización de saldo de envases retornables & Conductor repartidor & El sistema debe actualizar el saldo del cliente y del parque retornable en linea o en modo diferido según conectividad, asociando el registro a la entrega o a la guia correspondiente & Alta & Cap. 4.8, Cap. 9.7, Criterio Aceptación \#12 \\
\textbf{RF-06.11} & Confirmación de recepción mediante QR & Conductor repartidor & El sistema debe ofrecer al cliente la opción de confirmar la recepción mediante un código QR único generado por la app del conductor, cuando el cliente disponga de telefono inteligente. & Media & Cap. 16.1 (Decisión \#15) \\
\textbf{RF-06.12} & Confirmacion de recepción sin conectividad & Cliente , Conductor repartidor & El sistema debe permitir confirmar la recepción del cliente aun si este no dispone de telefono o conectividad, mediante captura de firma en pantalla o el registro fotografico de la recepcion, dejando evidencia vinculada a la guia de despacho electrónica en su acuse de recibo. & Media & Cap. 16.1 (Decisión \#15) \\
\textbf{RF-06.13} & Impresión de comprobante térmico en terreno. & Conductor repartidor propio y Cliente canal tradicional & La app de reparto debe permitir imprimir en terreno un comprobante térmico de entrega para el cliente usando una impresora portátil Bluetooth en aquellos casos donde el cliente requiera respaldo físico. & Media & RT-17.06 (impresora térmica portátil) \\
\textbf{RF-07.02} & Rendición Digital Individual de Efectivo & Cajero Liquidador / Conductor repartidor propio, Conductor repartidor externo & El sistema debe generar la planilla digital de rendición individual al retorno del camión, contrastando los valores declarados por el conductor en la aplicación móvil contra las entregas al contado realizadas y facturas cobradas. &  & Cap. 4.7, CA-10 \\
\textbf{RF-07.03} & Registro Obligatorio de Causales de Descuadre en Caja & Cajero Liquidador/ Conductor repartidor propio, Conductor repartidor externo & El sistema debe obligar al registro estructurado y categorizado de causales para cualquier diferencia detectada en la rendición de caja (descuentos comerciales autorizados, diferencias de cambio, siniestro en ruta), bloqueando el cierre de la liquidación si existe saldo sin justificar. &  & Cap. 4.7, 9.7, CA-10 \\
\textbf{RF-07.04} & Consola de Gestión de Cartera de Crédito & Ejecutivo de Cobranza/ Gerente de Finanzas, Clientes del canal tradicional, Clientes del canal moderno & El sistema debe proveer una consola de gestión para la cartera de crédito a 30 días sincronizada con el ERP, asignando carteras de clientes a ejecutivos de cobranza y registrando bitácoras de gestión, compromisos de pago y fechas de recontacto. &  & Cap. 4.7 \\
\textbf{RF-07.05} & Alerta y Bloqueo de Crédito en Preventa & Preventista / Ejecutivo de Cobranza, Gerente Comercial & El sistema debe notificar en la aplicación móvil de preventa el estado de bloqueo o mora de un cliente (>30 días o cupo excedido), bloqueando preventivamente el ingreso de pedidos a crédito y permitiendo únicamente venta al contado o cobro de deuda. &  & Cap. 4.3, 8, RT-09.01 \\
\textbf{RF-07.06} & Cobro con Terminal POS Móvil en Ruta & Conductor repartidor propio, Conductor repartidor externo / Clientes del canal tradicional, Clientes del canal moderno, Cajero Liquidador & El sistema debe integrarse vía Bluetooth con terminales de pago portátil (POS) en la aplicación móvil de reparto para cobro con tarjeta en ruta, capturando el código de autorización bancario y conciliando el voucher automáticamente en la liquidación diaria. &  & RT-17.06, Cap. 9.7 \\
\textbf{RF-07.07} & Consulta de Estado de Cuenta en Portal de Clientes & Clientes del canal tradicional, Clientes del canal moderno/ Ejecutivo de Cobranza, Gerente de Finanzas & El sistema debe disponer en el portal web de autoatención el estado de cuenta y la cartola de antigüedad de deuda (tramos corriente, 30, 60, 90+ días) descargable en formatos PDF y Excel para clientes autenticados. &  & Cap. 9.7, RT-16.30 \\
\textbf{RF-07.08} & Validación en Tesorería de Pagos Web & Cajero Liquidador / Clientes del canal moderno, Gerente de Finanzas & El sistema debe proveer una bandeja de validación para pagos reportados en el portal web de clientes, permitiendo al cajero verificar el comprobante bancario y autorizar su traspaso estructurado al ERP. &  & Cap. 4.7, 9.7, RT-16.30 \\
\textbf{RF-07.09} & Interfaz Asíncrona de Cobranzas hacia ERP & Jefe de TI / Gerente de Finanzas, Cajero Liquidador & El sistema debe transferir de forma estructurada, asíncrona e idempotente los comprobantes de recaudación y cierres comerciales aprobados hacia el ERP para actualizar automáticamente las cuentas por cobrar. &  & Cap. 5, Restricción R4 \\
\textbf{RF-07.10} & Registro de Abonos y Pagos Parciales en Entrega & Conductor repartidor propio, Conductor repartidor externo / Clientes del canal tradicional, Cajero Liquidador, Gerente de Finanzas & El sistema debe permitir el registro de anticipos o pagos parciales en entregas al contado desde la aplicación móvil de reparto, emitiendo un comprobante digital de abono y traspasando el saldo insoluto a la cuenta corriente del cliente en el ERP. &  & Cap. 4.7 \\
\textbf{RF-07.11} & Cobranza Presencial de Facturas por Preventistas & Preventista / Cajero Liquidador, Clientes del canal tradicional & El sistema debe permitir al preventista registrar la cobranza en efectivo de facturas vencidas en el local del cliente desde su terminal móvil, emitiendo un recibo provisional digital y forzando su rendición diaria en caja. &  & Cap. 4.3, Cap. 8 \\
\textbf{RF-08.01} & Registro de Devoluciones en Ruta & Conductor repartidor & El sistema debe permitir al conductor registrar la devolución de productos en el momento de la entrega, capturando SKU, cantidad y motivo (daño / vencimiento / producto no pedido / error), operando en modo offline. & Muy Alta & Cap. 4.6, Cap. 4.8, RT-17.01 \\
\textbf{RF-08.02} & Gestión del Destino de Productos Devueltos & Jefe de bodega & Al recepcionar una devolución en el andén del CD, el sistema debe permitir seleccionar el destino final del producto (reingreso a stock / venta con descuento / devolución al proveedor / destrucción) y enviar la transacción al ERP mediante su interfaz de integración. & Muy Alta & Cap. 4.8, Cap 10 (restriccion 4) \\
\textbf{RF-08.03} & Gestión de Cuenta Corriente de Envases Retornables & Conductor repartidor & El sistema debe mantener una cuenta corriente de envases retornables por cliente, actualizando el saldo en cada entrega y devolución registrada por el conductor. & Alta & Cap. 4.8, Cap. 16.1 (Decisión \#10), Cap 9.7 \\
\textbf{RF-08.04} & Alertas por Exceso de Envases Retornables & Preventista & El sistema debe generar una alerta al preventista cuando el saldo de envases retornables de un cliente supere el umbral configurado por el administrador. & Alta & Cap. 4.8, Cap 9.7 \\
\textbf{RF-08.05} & Registro de Mermas por Vencimiento & Preparador de bodega / Preventista & El sistema debe permitir registrar mermas por vencimiento detectadas en conteo cíclico, preparación de pedidos o góndola del cliente, asociando cada registro al SKU y al lote correspondiente. & Alta & Cap. 4.8, Cap 10 (restriccion 4) \\
\textbf{RF-08.06} & Reporte de Pérdida de Envases Retornables & Jefe de operaciones / Finanzas & El sistema debe reportar la pérdida de envases retornables desglosada por cliente, tipo de envase y zona, con actualización diaria. & Alta & Cap. 4.8, Cap 18 (Criterio Aceptación \#12) \\
\textbf{RF-08.07} & Integración de Devoluciones y Mermas al Costo de Servir & Finanzas & El sistema debe incluir el costo de devoluciones y mermas por cliente como componente del cálculo del costo de servir (Épica 11). & Alta & Cap. 9.8 \\
\textbf{RF-09.01} & Trazabilidad de producto foward tracing & Equipo de calidad & El sistema debe permitir la trazabilidad 'hacia adelante' (forward tracing): dado un lote específico, debe listar todos los clientes que recibieron productos de ese lote, con fecha de entrega, guía y cantidad. & Muy Alta & Cap. 1, Cap. 9.6, Criterio Aceptación \#1 \\
\textbf{RF-09.02} & Trazabilidad de producto backward tracing & Equipo de calidad & El sistema debe permitir la trazabilidad 'hacia atrás' (backward tracing): dado un producto en el local del cliente (SKU + lote), debe identificar el origen (proveedor, fecha de recepción y orden de compra) & Muy Alta & Cap. 4.9, Cap. 9.6 \\
\textbf{RF-09.03} & Registro de temperaturas de cámaras de refigerado y congelado & Sensor de temperatura & El sistema debe registrar de forma automática y continua la temperatura de las cámaras de refrigerado y congelado, con lecturas cada 5 minutos como mínimo & Muy Alta & Cap. 4.9, Cap. 7.3, Criterio Aceptación \#3 \\
\textbf{RF-09.04} & Registro de temperaturas de camiones con equipo de frío & Sensor de temperatura & El sistema debe registrar de forma automática y continua la temperatura en los camiones con equipo de frío (propios + de transportistas que cuenten con sensor), durante todo el trayecto desde la carga hasta la entrega. & Muy Alta & Cap. 4.6, Cap. 4.9, Criterio Aceptación \#3 \\
\textbf{RF-09.05} & Alerta de excusiones térmicas & Sistema de monitoreo & El sistema debe generar una alerta en tiempo real (notificación en cabina y en plataforma de monitoreo) cuando se detecte una excursión de temperatura que supere el umbral definido para el tipo de producto. & Muy Alta & Cap. 8 (Katherine), Cap. 16.1 (Decisión \#4) \\
\textbf{RF-09.06} & Parametrización de excusiones térmicas & Jefe de calidad/operaciones & El sistema debe permitir configurar los parámetros que determine una excursión térmica (definida por tiempo/grados) & Muy Alta & Cap. 8 (Katherine vs Nelson), Cap. 16.1 (Decisión \#4) \\
\textbf{RF-09.07} & Bloqueo de despacho por excusiones térmicas & Conductor & El sistema permitira al conductor bloquear el despacho del cambion según su criterio, basandose en el tipo de producto y criterio del mandante & Media & Cap. 8 (Katherine vs Nelson), Cap. 16.1 (Decisión \#4) \\
\textbf{RF-09.08} & Registro de pedido rechazado por cliente & Conductor & El sistema debe permitir al conductor registrar el rechazo de un producto entregado por 'sospecha de ruptura de cadena de frío', asociando el evento al sensor de temperatura del camión en ese tramo. & Muy Alta & Cap. 4.9, Cap. 8 \\
\textbf{RF-09.09} & Registro de lotes con control sanitario & Preparador & El sistema debe registrar, en el momento de la preparación (picking), el lote específico que se le asigna a cada línea de pedido, para los productos con control sanitario. & Muy Alta & Cap. 4.5, Cap. 4.9 \\
\textbf{RF-09.10} & Generación de cumplimiento de cadena de frío &  & El sistema debe generar la evidencia de cumplimiento de cadena de frío exigible ante la autoridad sanitaria y ante el cliente, para un envío determinado. & Muy Alta & Cap. 4.9 ; Cap. 18 criterio N\ensuremath{^\circ}3 \\
\textbf{RF-11.01} & Distribución Automática de Reportes Ejecutivos & Gerente de Finanzas / Gerente Comercial, Jefa de Calidad & El sistema debe programar y distribuir vía correo electrónico reportes ejecutivos consolidados (en formatos PDF y XLSX) con periodicidad semanal y mensual, desglosando la evolución de OTIF, Fill Rate, Merma por vencimiento y Costo de Servir por zona y canal de comercialización. &  & Cap. 7.1, 7.2, 9.8, CA-11 \\
\textbf{RF-11.02} & Cálculo del Costo de Servir por Entrega y Cliente & Gerente de Finanzas / Gerente Comercial & El sistema debe calcular el Costo de Servir real por entrega y por cliente, imputando los costos directos de transporte (combustible devengado vía GPS, peajes, tiempo de conducción y tiempo de descarga) y los costos indirectos prorrateados (bodega y administración). &  & Cap. 4.4, 7.2, 9.8, CA-11 \\
\textbf{RF-11.03} & Cálculo Diario Unificado del Indicador OTIF & Gerente Comercial (Primario) / Jefa de Calidad, Conductor repartidor propio, Conductor repartidor externo (Secundarios) & El sistema debe calcular diariamente el indicador OTIF por cliente, ruta, zona y total consolidado, definiendo «On-Time» como el cumplimiento de fecha y ventana horaria pactada e «In-Full» como la entrega del 100\% de ítems y unidades pedidas, obligando al registro de causal tipificada ante desvíos. &  & Cap. 7.1, 9.1, CA-04 \\
\textbf{RF-11.04} & Medición de Fill Rate y Perfect Order con POD Móvil & Gerente Comercial (Primario) / Conductor repartidor propio, Conductor repartidor externo (Secundarios) & El sistema debe calcular automáticamente el Fill Rate (unidades entregadas vs. pedidas) y el índice de Perfect Order contrastando la orden de venta original contra los datos de entrega física confirmados en el dispositivo móvil de reparto, sin requerir interacción tecnológica del cliente tradicional. &  & Cap. 7.1, 9.4, 16.2, CA-08 \\
\textbf{RF-11.05} & Monitoreo y Alerta de Ocupación de Flota & Planificador de rutas / Gerente Comercial & El sistema debe calcular diariamente la utilización de capacidad volumétrica (m\textasciicircum{}3) y de peso (kg) por viaje despachado, alertando aquellos camiones que salgan con una ocupación inferior al umbral configurable (ej. < 65\%). &  & Cap. 4.4, 7.2, CA-14 \\
\textbf{RF-11.06} & Tablero de Control Operacional en Tiempo Real & Gerente Comercial Gerente de Finanzas, Jefa de Calidad & El sistema debe proveer un Tablero de Control Operacional en tiempo real que consolide el avance de rutas de despacho, porcentaje de entregas cumplidas, devoluciones en tránsito, alertas activas de cadena de frío y monto recaudado en efectivo. &  & RT-05.29, Cap. 9.1, 9.6 \\
\textbf{RF-11.07} & Liquidación y Cierre Comercial por Camión & Cajero Liquidador / Conductor repartidor propio, Conductor repartidor externo, Gerente de Finanzas & El sistema debe procesar automáticamente la liquidación y cierre comercial individual por camión al retornar a base, consolidando documentos tributarios entregados, devoluciones físicas recibidas, cobranzas en efectivo y balance de envases retornables. &  & RT-05.29, CA-10 \\
\textbf{RF-11.08} & Segmentación de Clientes por Rentabilidad Neta & Gerente Comercial / Preventista, Gerente de Finanzas & El sistema debe segmentar mensualmente la cartera de clientes en matrices de rentabilidad neta (margen comercial vs. costo de servir real), sugiriendo ajustes de frecuencia de visita o umbrales mínimos de compra para clientes no rentables. &  & Cap. 7.2, 9.8, Decisión 16.1 \#7 \\
\textbf{RF-12.01} & Recepción de pedido electrónico & Cliente del canal moderno (Cadena de supermercados) & El sistema debe recibir el mensaje de pedido electrónico estructurado (EDI/EDIFACT/XML, según el estándar que exija cada cadena registrada del canal moderno). & Alta & Cap. 9.9; RT-05.23 \\
\textbf{RF-12.02} & Validación automática del pedido EDI & Cliente del canal moderno (Cadena de supermercados) & El sistema debe validar automáticamente el pedido EDI recibido, verificando estructura del mensaje, existencia del cliente, existencia de los productos, cantidades, precios y demás datos obligatorios de la cadena. & Alta & Cap. 9.9 / Cap. 17.1 \\
\textbf{RF-12.03} & Resolución del código de producto contra el maestro de Puelche & Cliente del canal moderno (Cadena de supermercados) & El sistema debe resolver el código de producto informado por la cadena contra el maestro de productos vigente de Puelche, mediante una tabla de equivalencias por cadena. & Alta & Cap. 16.1, Decisión \#11 \\
\textbf{RF-12.04} & Registro automático del pedido EDI válido en el sistema de gestión & Cliente del canal moderno (Cadena de supermercados) & El sistema debe registrar automáticamente, sin digitación manual, el pedido EDI válido (RF-12.02, RF-12.03) en el sistema de gestión empresarial. &  &  \\
\textbf{Criterio de Aceptación \#13} &  &  &  &  &  \\
\textbf{RF-12.05} & Disponibilización del pedido a la planificación de despacho & Planificador de rutas & El sistema debe poner a disposición del proceso de planificación de rutas todo pedido registrado proveniente de un canal electrónico. & Alta & Cap. 18, Criterio de Aceptación \#13 \\
\textbf{RF-12.06} & Gestión de excepciones de pedidos EDI inválidos & Operador & El sistema debe permitir gestionar, mediante una bandeja de excepciones, los pedidos que no pasaron la validación (RF-12.02) o la resolución de maestro (RF-12.03), permitiendo corregir, rechazar o reprocesar cada uno. & Alta & Cap. 17.1 / 17.2 \\
\textbf{RF-12.07} & Confirmación o rechazo del pedido a la cadena de origen & Cliente del canal moderno (Cadena de supermercados) & El sistema debe enviar a la cadena de origen una confirmación o rechazo del pedido, indicando la causa cuando no pueda ser aceptado completamente. & Alta & Cap. 9.9 \\
\textbf{RF-12.08} & Historial de estados del pedido electrónico & Operador de excepciones EDI / Planificador de rutas & El sistema debe mantener el historial de estados de cada pedido electrónico (recepción, validación, confirmación, preparación, despacho, ASN y prueba de entrega) &  &  \\
\textbf{hora y responsable/sistema,Alta,Cap. 9.1 / 9.4 / 9.9} &  &  &  &  &  \\
\textbf{RF-12.09} & Envío de aviso de despacho anticipado (ASN) & Cliente del canal moderno (Cadena de supermercados) & El sistema debe generar y enviar al cliente del canal moderno un ASN cuando su pedido queda cargado en el camión, referenciando la guía de despacho electrónica ya emitida, con productos, cantidades y hora estimada de llegada. & Alta & Cap. 9.9; entrevista Rubén Salinas \\
\textbf{RF-12.10} & Planificación de la llegada dentro de la ventana de 30 minutos & Planificador de rutas & El sistema debe planificar la llegada del camión al cliente del canal moderno dentro de la ventana horaria de 30 minutos comprometida. & Alta & Cap. 9.9 \\
\textbf{RF-12.11} & Registro del cumplimiento de la ventana de entrega & Conductor & El sistema debe registrar el cumplimiento o incumplimiento de la ventana de 30 minutos para cada entrega del canal moderno, distinguiéndola de la ventana crítica de despacho 05:30-07:00. &  &  \\
\textbf{RF-12.12} & Captura de evidencia de entrega en canal moderno & Conductor & El sistema debe capturar evidencia de entrega en el local del cliente: firma en pantalla del titular cuando pueda firmar, o fotografía de la mercadería más nombre de quien recibe cuando no sea el titular o no pueda firmar. & Alta & Cap. 9.9; Cap. 16.1, Decisión \#15 \\
\textbf{RF-12.13} & Envío de acuse de recibo digital al cliente & Cliente canal moderno & El sistema debe enviar al cliente del canal moderno un acuse de recibo digital de la entrega, articulado con la guía de despacho electrónica y su acuse de recibo tributario. & Alta & Cap. 9.9 \\
\textbf{RF-12.14} & Portal público de catálogo & Usuario no autenticado & El sistema debe publicar un portal web público, sin autenticación, con el catálogo vigente de productos (sin precios ni stock), contacto comercial e información institucional. & Media & RT-16.30 \\
\textbf{RF-12.15} & Pedido en autoservicio para clientes del canal moderno & Cliente & El portal autenticado de clientes debe permitir generar un pedido en autoservicio, validando el mismo crédito y stock que consulta el preventista. & Alta & RT-16.30 \\
\textbf{RF-12.16} & Consulta de estado de entregas en el portal de clientes & Cliente & El portal de clientes debe permitir consultar el estado de sus entregas y la ventana de entrega comprometida. & Alta & RT-16.30 / Cap. 9.1 \\
\textbf{RF-12.17} & Consulta de documentos tributarios en el portal de clientes & Cliente & El portal de clientes debe permitir consultar y descargar sus documentos tributarios (factura, guía electrónica, nota de crédito). & Alta & RT-16.30 / RT-16.14 \\
\textbf{RF-12.18} & Consulta de saldo y estado de cuenta en el portal de clientes & Cliente & El portal de clientes debe permitir consultar el saldo y estado de cuenta del cliente autenticado, con datos del sistema de gestión. & Alta & RT-16.30 \\
\textbf{RF-12.19} & Consulta de rutas asignadas en el portal de transportistas & Empresa transportista / Jefe de flota & El portal de transportistas debe permitir consultar las rutas asignadas a la empresa para el día siguiente. & Media & RT-16.30 \\
\textbf{RF-12.20} & Descarga de documentación de despacho en el portal de transportistas & Empresa transportista / Jefe de flota & El portal de transportistas debe permitir descargar la documentación de despacho de cada ruta asignada. &  &  \\
\textbf{RF-12.21} & Confirmación diaria de conductor y vehículo & Empresa transportista / Jefe de flota & El portal debe permitir confirmar, antes de la ventana de despacho, qué conductor y vehículo ejecutarán cada ruta asignada del día. & Media & Cap. 16.1, Decisión \#5 \\
\textbf{RF-12.22} & Consulta de órdenes de compra en el portal de proveedores & Proveedor & El portal de proveedores debe permitir consultar, en modo de solo lectura, el estado de sus órdenes de compra. & Media & RT-16.30 \\
\textbf{RF-12.23} & Consulta de recepciones en el portal de proveedores & Proveedor & El portal de proveedores debe permitir consultar, en modo de solo lectura, las recepciones registradas para sus órdenes. & Media & RT-16.30 \\
\textbf{RF-12.24} & Consulta de devoluciones y notas de crédito en el portal de proveedores & Proveedor & El portal de proveedores debe permitir consultar, en modo de solo lectura, sus devoluciones y notas de crédito. &  &  \\
\textbf{RF-12.25} & Registro de cliente en el portal & Cliente & El portal debe permitir al cliente del canal moderno registrarse con su RUT de empresa, correo electrónico corporativo y una contraseña, siempre que el RUT esté previamente habilitado en el sistema de gestión como cliente activo de Puelche. & Muy Alta &  \\
\textbf{RF-12.26} & Inicio de sesión en el portal de clientes & Cliente & El portal debe permitir al cliente autenticarse con su correo y contraseña registrados. & Muy Alta &  \\
\textbf{RF-12.27} & Recuperación de contraseña & Cliente & El portal debe permitir al cliente solicitar un enlace de restablecimiento de contraseña a su correo registrado cuando no recuerde su clave. & Alta &  \\
\textbf{RF-12.28} & Cierre de sesión en el portal de clientes & Cliente & El portal debe permitir al cliente cerrar su sesión de forma explícita, invalidando el token de sesión activo. & Alta &  \\
\textbf{RF-12.29} & Expiración automática de sesión por inactividad & Cliente & El portal debe cerrar automáticamente la sesión del cliente tras un período de inactividad definido, sin requerir acción explícita de cierre. & Alta &  \\
\textbf{RF-14.01} & Integración con telemetría de camiones propios & Jefe de operaciones / Gerente de operaciones & El sistema debe integrarse con la plataforma de telemetría de terceros existente que cubre los camiones propios para consumir datos de posición GPS, velocidad y kilometraje. & Muy Alta & Cap. 5, RT-17.06 \\
\textbf{RF-14.02} & Extensión de telemetrá a camiones externos & Jefe de flota / Transportista, Jefe de operaciones & El sistema debe extender la visibilidad de telemetría a los camiones de transportistas externos contratados cuyo operador cuente con GPS habilitado, mediante un mecanismo acordado mediante la aplicacion. & Alta & Cap. 2.3, Cap. 17.4 (punto 9). \\
\textbf{RF-14.03} & Visualización de rutas planificadas vs reales & Gerente de operaciones, Planificador de rutas & El sistema debe mostrar en el mapa las rutas planificadas vs las rutas reales recorridas por cada camión, para identificar desviaciones y optimizar las rutas futuras. & Alta & Cap. 3 y entrevista de nelson \\
\textbf{RF-14.04} & Alerta de desviación de ruta & Centro de control / Jefe de operaciones & El sistema debe generar una alerta cuando el camión se desbie más de 3 km de la ruta planificada, para que el centro de control pueda reaccionar. & Alta & Cap. 9.1 \\
\textbf{RF-14.05} &  & Conductor repartidor, Jefe de flota & El sistema debe registrar el tiempo de conducción y descanso del conductor, utilizando los datos de posicionamiento (con la debida consideración de la objeción sindical) para cumplir con el régimen especial de jornada. & Media & Cap. 12 (Régimen jornada conductores), Cap. 16.1 (Decisión \#13), Restricción \#10 \\
\textbf{RF-14.06} & Integración de telemetría con costo de servir & Gerente de finanzas, Jefe de operaciones & El sistema debe integrar los datos de telemetría con el módulo de costo de servir, imputando kilómetros recorridos y tiempo de viaje a cada entrega y cliente. & Muy Alta & Cap. 9.8 \\
\textbf{RF-14.07} & Vinculación conductor-camión en cada viaje & Conductor repartidor, Jefe de flota & El sistema debe registrar qué conductor propio o de transportista externo está asociado a cada camión en cada viaje, consumiendo la identidad validada por el mecanismo de autenticación (RF-15.02). Esta vinculación es necesaria para imputar responsabilidades en cobros, entregas y eventos de ruta, y para resolver la Decisión \#5. & Alta & Cap. 16.1 (Decisión \#5), Cap. 4.6 \\
\textbf{RF-14.08} & Gestión de geocercas para control de llegada/salida & Jefe de operaciones / Gerente de operaciones, Conductor repartidor & El sistema debe permitir la gestión de geocercas zonas geográficas, para identificar cuándo un camión entra o sale de una zona ej: CD Talca, zona de clientes, para medir tiempos de llegada y salida. & Media & Cap. 9.9 y entrevista con ruben salinas \\
\textbf{RF-14.09} & Reportes de kilometraje y consumo de combustible & Gerente de finanzas, Jefe de operaciones & El sistema debe consolidar los datos de telemetría de todos camiones propios para generar reportes de kilometraje total y consumo de combustible estimado, apoyando la gestión de costos de transporte. & Alta & Cap. 7.2 (costo desconocido). \\
\end{longtable}



\subsection{Requerimientos no funcionales (RNF)}
\begin{longtable}{p{1.8cm} p{2.3cm} p{1.8cm} p{4.8cm} p{2.2cm} p{1.1cm}}
\caption{Cat\'alogo de requerimientos no funcionales (RNF)}\label{tab:rnf}\\
\toprule
\textbf{ID} & \textbf{T\'itulo} & \textbf{Actor} & \textbf{Descripci\'on} & \textbf{Precondici\'on} & \textbf{Prioridad} \\
\midrule
\endfirsthead
\toprule
\textbf{ID} & \textbf{T\'itulo} & \textbf{Actor} & \textbf{Descripci\'on} & \textbf{Precondici\'on} & \textbf{Prioridad} \\
\midrule
\endhead
\bottomrule
\endlastfoot
\textbf{RNF-12.01} & Hito de producción del canal moderno & Equipo de proyecto & El portal de autoatención de clientes (RF-12.15 a RF-12.18) y la recepción EDI (RF-12.01 a RF-12.09) deben estar en producción antes de enero de 2029. & Ninguna. & Muy Alta \\
\textbf{RNF-12.02} & Configurabilidad de integración EDI para nuevas cadenas & ---- & El mecanismo de recepción EDI debe permitir configurar el mapeo de mensajería de una nueva cadena sin desarrollo de software a medida. & Existe al menos una integración EDI ya operativa. & Alta \\
\textbf{RNF-03-01} & Tiempo de respuesta de registro de línea de pedido & ---- & El tiempo para registrar una línea de pedido no debe superar 1,5 segundos, con o sin señal parcial. & ---- & Muy Alta \\
\textbf{RNF-03-02} & Tiempo de respuesta de consulta de stock y crédito con conectividad & ---- & El tiempo para consultar stock y crédito con conectividad no debe superar 2 segundos. & ---- & Muy Alta \\
\textbf{RNF-03-03} & Tiempo de respuesta de consulta de historial de compra & ---- & El tiempo para consultar el historial de compra con conectividad no debe superar 2 segundos. & ---- & Muy Alta \\
\textbf{RNF-01.01} &  &  & El tiempo de procesamiento para el registro de cada ítem en recepción, desde la confirmación del escaneo hasta la respuesta en pantalla, no debe superar los 2 segundos. &  & Alta \\
\textbf{RNF-01.02} &  &  & El sistema debe conservar los registros de recepción según el plazo por tipo de dato: trazabilidad sanitaria y no conformidades por mínimo 5 años; documentos de referencia tributaria por mínimo 6 años. &  & Muy Alta \\
\textbf{RNF-01.03} &  &  & El sistema debe operar en modo desconectado durante mínimo 24 horas continuas sin pérdida de datos y sincronizar todas las transacciones pendientes con el ERP en $\le$ 2 horas tras restablecer la conectividad. &  & Muy Alta \\
\textbf{RNF-02.01} & Operación offline del CD (24 horas) & — & El sistema de gestión de almacén (inventario, movimientos, conteos, consultas de stock) debe operar en modo desconectado durante un mínimo de 24 horas continuas sin pérdida de datos. Al restablecer la conectividad, las transacciones pendientes deben sincronizarse con el sistema central en $\le$ 2 horas, con resolución determinista de conflictos de stock (ej. si dos operarios movieron el mismo pallet). & El CD opera con conectividad normal y se produce un corte de enlace. & Muy Alta \\
\textbf{RNF-05.01} & Latencia de confirmación de picking ($\le$ 1s) & — & La confirmación de una línea de picking (escaneo + validación + almacenamiento local) no debe superar 1 segundo de tiempo de respuesta en el dispositivo del preparador, tanto en modo online como en modo offline, incluyendo la operación dentro de la cámara de congelado ($-$22 \ensuremath{^\circ}C). & El preparador está ejecutando una misión de picking. & Muy Alta \\
\textbf{RNF-05.02} & Operabilidad con guantes a $-$22 \ensuremath{^\circ}C & — & La interfaz de picking debe ser operable sin precisión de toque fino, con elementos interactivos de al menos 15$\times$15 mm y sin requerir gestos complejos (deslizar, pellizcar, doble toque). El dispositivo debe mantener funcionalidad operativa durante al menos 30 minutos continuos en cámara de congelado a $-$22 \ensuremath{^\circ}C. & El preparador opera en la cámara de congelado con guantes térmicos. & Muy Alta \\
\textbf{RNF-05.03} & Curva de aprendizaje ($\le$ 2 horas) & — & Un nuevo operario de preparación, sin experiencia previa con el sistema, debe completar una misión de picking de al menos 20 líneas con una tasa de error $\le$ 5\% tras un máximo de 2 horas de entrenamiento asistido en el puesto de trabajo. La capacitación no debe requerir materiales impresos de más de 2 páginas ni sesiones fuera del recinto de bodega. & Nuevo operario ingresa al turno de preparación. & Muy Alta \\
\textbf{RNF-08.01} &  &  & El registro de una devolución en el andén del CD (escaneo + selección de destino) no debe superar 30 segundos por ítem. &  & Alta \\
\textbf{RNF-08.02} &  &  & Los registros de devoluciones y mermas deben conservarse según el plazo por tipo de dato: trazabilidad sanitaria por mínimo 5 años; documentos con efectos tributarios (notas de crédito asociadas) por mínimo 6 años. &  & Alta \\
\textbf{RNF-09.01} &  &  & El sistema tiene un tiempo de respuesta para generar el listado de clientes afectados por un lote (RF-09.02) que no debe superar las 2 horas (incluyendo la generación del archivo exportable). &  &  \\
\textbf{RNF-09.02} &  &  & El sistema debe conservar los registros sanitarios y de trazabilidad de lote durante el período mínimo legal: 'vida útil del producto + 6 meses', con un piso mínimo de 5 años. &  &  \\
\textbf{RNF-09.03} &  &  & El sistema debe registrar todas las consultas a datos de trazabilidad (lotes, clientes afectados) y modificaciones al maestro de lotes en un log de auditoría, con identificación del usuario y fecha/hora. &  &  \\
\textbf{RNF-09.04} &  &  & El sistema debe funcionar Sin cobertura de señal y sincronizando al recuperar conectividad dentro de la cámara de frío &  &  \\
\textbf{RNF-07.01} &  & Jefe de TI / Cajero Liquidador, Gerente de Finanzas & El sistema debe sincronizar la información de preconciliación de cobranzas capturada en la aplicación móvil en un tiempo no superior a 10 minutos tras conectar a red, permitiendo la liquidación de la flota el mismo día. &  & Muy Alta \\
\textbf{RNF-07.02} &  &  & El sistema debe cifrar a nivel de campo (AES-256 en reposo, TLS 1.3 en tránsito) los datos de pagos en efectivo, transacciones con tarjeta POS y comportamiento comercial de clientes (Ley N\ensuremath{^\circ} 21.719). &  & Muy Alta \\
\textbf{RNF-04.01} &  &  & El motor de ruteo debe procesar la asignación y secuencia completa de la matriz diaria de entregas en un tiempo acotado. &  & Media \\
\textbf{RNF-06.01} &  &  & El sistema debe persistir el 100\% de los datos necesarios en el almacenamiento cifrado local para ejecutar la ruta sin cobertura. &  & Media \\
\textbf{RNF-11.01} &  &  & El sistema debe desplegar los indicadores operacionales del día con una latencia no superior a 5 minutos desde su recepción en el backend central para terminales con cobertura móvil activa. &  & Muy Alta \\
\textbf{RNF-11.02} &  &  & El sistema debe desacoplar el procesamiento analítico y de reportería (OLAP/BI) del motor transaccional operativo (OLTP), garantizando que las consultas analíticas pesadas no degraden los tiempos de respuesta de transacciones críticas (picking \$\textbackslash\{\}le\$ 1s, preventa \$\textbackslash\{\}le\$ 1,5s). &  & Alta \\
\end{longtable}



\subsection{Requerimientos de bases (RT)}

\subsubsection{Requerimientos funcionales de bases}
\begin{longtable}{p{1.7cm} p{2.3cm} p{1.9cm} p{5.0cm} p{1.1cm} p{1.9cm}}
\caption{Requerimientos funcionales de Bases (RT)}\label{tab:bt_rf}\\
\toprule
\textbf{ID} & \textbf{T\'itulo} & \textbf{Actor} & \textbf{Descripci\'on} & \textbf{Prioridad} & \textbf{Origen} \\
\midrule
\endfirsthead
\toprule
\textbf{ID} & \textbf{T\'itulo} & \textbf{Actor} & \textbf{Descripci\'on} & \textbf{Prioridad} & \textbf{Origen} \\
\midrule
\endhead
\bottomrule
\endlastfoot
\textbf{RF-13.01} & Declaración de arquitectura multicapa & Equipo de proyecto / Arquitecto & El PROPONENTE debe declarar la arquitectura lógica de la solución organizada en las 8 capas del modelo de referencia: Presentación, Borde y Exposición, Puerta de Enlace, Servicios de Negocio, Integración y Eventos, Datos, Seguridad Transversal y Observabilidad Transversal. & Muy Alta & RT-02.01, RT-02.03 \\
\textbf{RF-13.02} & Tabla de emplazamiento nube/on-premise & Equipo de proyecto & El PROPONENTE debe presentar una tabla de emplazamiento que asigne cada componente de la solución a nube o a on-premise, justificando cada decisión según el criterio de emplazamiento (latencia, criticidad, volumen de datos, acoplamiento físico, elasticidad, regulación). & Muy Alta & Cap. 3.1, RT-03.01 \\
\textbf{RF-13.03} & Registro de decisiones de arquitectura (ADR) & Equipo de proyecto & El PROPONENTE debe mantener un registro de decisiones de arquitectura (ADR) fechado, con la alternativa escogida, las alternativas descartadas y el criterio de decisión, actualizado durante toda la ejecución. & Alta & RT-02.04 \\
\textbf{RF-13.04} & Modelo de dominio de negocio & Equipo de proyecto & El PROPONENTE debe presentar un modelo de dominio del negocio del caso, con las entidades principales, sus relaciones y los eventos de negocio que las modifican. & Muy Alta & RT-02.13 \\
\textbf{RF-14.01} & Clasificación de información por nivel de sensibilidad & Administrador de seguridad & El sistema debe permitir clasificar la información del CLIENTE por nivel de sensibilidad (pública, interna, confidencial, restringida), con controles diferenciados por nivel documentados en una matriz. & Muy Alta & RT-11.03 \\
\textbf{RF-14.02} & Matriz de controles de seguridad trazable a ISO 27001/27002 & Administrador de seguridad & El PROPONENTE debe mantener una matriz de controles de seguridad trazable a ISO/IEC 27001 e ISO/IEC 27002, indicando el control, su implementación concreta en la solución y la evidencia que lo acredita. & Muy Alta & RT-11.05 \\
\textbf{RF-14.03} & Declaración de superficie de exposición & Administrador de seguridad & El PROPONENTE debe declarar la superficie de exposición completa de la solución: cada nombre de dominio, puerto y servicio alcanzable desde fuera de la red del CLIENTE. & Alta & RT-11.13 \\
\textbf{RF-14.04} & Plan de respuesta a incidentes de seguridad & Administrador de seguridad & El PROPONENTE debe disponer de un plan de respuesta a incidentes de seguridad que defina clasificación, cadena de escalamiento, plazos, responsables y protocolo de comunicación al CLIENTE dentro de las 2 horas de detectado un incidente de severidad crítica. & Muy Alta & RT-11.18 \\
\textbf{RF-14.05} & Notificación de brecha de seguridad al CLIENTE & Administrador de seguridad & El sistema (o el equipo de seguridad del ADJUDICATARIO) debe notificar al CLIENTE cualquier brecha de seguridad o de datos personales dentro de las 24 horas de su detección, con informe preliminar, y entregar el análisis de causa raíz dentro de los 5 días hábiles siguientes. & Muy Alta & RT-11.19 \\
\textbf{RF-15.01} & Gestión centralizada de identidad & Administrador del sistema & El sistema debe contar con una gestión de identidad centralizada, con federación mediante OpenID Connect y OAuth 2.1, o SAML 2.0 cuando la integración con el CLIENTE lo requiera, e integración con el directorio corporativo del CLIENTE por LDAP o equivalente en la nube. & Muy Alta & RT-12.01 \\
\textbf{RF-15.02} & Inicio de sesión único (SSO) para todos los módulos & Administrador del sistema & El sistema debe implementar inicio de sesión único para todos los módulos de la solución, con cierre de sesión propagado a todos ellos. & Muy Alta & RT-12.02 \\
\textbf{RF-15.03} & Autenticación multifactor (MFA) para administradores y acceso externo & Administrador del sistema & El sistema debe exigir autenticación multifactor para personas usuarias administradoras, para todo acceso privilegiado y para todo acceso originado fuera de la red corporativa, soportando factores FIDO2 o claves de acceso. & Muy Alta & RT-12.03 \\
\textbf{RF-15.04} & Control de acceso basado en roles (RBAC) y segregación de funciones & Administrador del sistema & El sistema debe implementar control de acceso basado en roles, complementado con control basado en atributos donde el proceso lo exija, con matriz de segregación de funciones documentada y verificable. & Muy Alta & RT-12.05 \\
\textbf{RF-15.05} & Aprovisionamiento y desaprovisionamiento automatizado de cuentas & Administrador del sistema & El sistema debe automatizar el aprovisionamiento y desaprovisionamiento de cuentas, ligado al ciclo de vida laboral, con baja efectiva en un plazo no superior a 24 horas desde la desvinculación. & Alta & RT-12.10 \\
\textbf{RF-16.01} & Observabilidad unificada para nube y on-premise & Administrador del sistema / Equipo de operaciones & El sistema debe implementar observabilidad unificada para nube y on-premise, con métricas, registros y trazas distribuidas correlacionadas por un identificador único de transacción, instrumentadas conforme a OpenTelemetry. & Muy Alta & RT-14.01 \\
\textbf{RF-16.02} & Acceso del CLIENTE a tableros operacionales y de negocio & Administrador del sistema / CLIENTE & El CLIENTE debe disponer de acceso propio y permanente a los tableros operacionales y de negocio, con datos en tiempo real y capacidad de exportación. & Muy Alta & RT-14.02 \\
\textbf{RF-16.03} & Alertamiento basado en síntomas de negocio & Administrador del sistema / Equipo de operaciones & El sistema debe implementar alertamiento basado en síntomas de negocio y no sólo en umbrales de infraestructura, con supresión de ruido, agrupación, escalamiento automático y turnos de disponibilidad declarados. & Alta & RT-14.04 \\
\textbf{RF-16.04} & Análisis de causa raíz obligatorio tras incidente crítico & Equipo de operaciones / Gerente de servicio & El sistema (o el equipo del ADJUDICATARIO) debe ejecutar un análisis de causa raíz obligatorio tras cada incidente crítico, con informe entregable dentro de cinco días hábiles y seguimiento de las acciones correctivas hasta su cierre. & Alta & RT-14.06 \\
\textbf{RF-17.01} & Módulo de administración de usuarios, roles y permisos & Administrador del sistema & El sistema debe disponer de un módulo de administración que permita al CLIENTE, sin intervención del ADJUDICATARIO, gestionar personas usuarias, roles, permisos, unidades organizacionales y sus jerarquías. & Muy Alta & RT-16.01 \\
\textbf{RF-17.02} & Parametrización de reglas de negocio desde interfaz de administración & Administrador del sistema & El sistema debe permitir parametrizar reglas de negocio (umbrales, plazos, montos, tolerancias, catálogos, listas de valores, textos de notificación) desde la interfaz de administración, con control de versiones y registro de quién cambió qué y cuándo. & Muy Alta & RT-16.02 \\
\textbf{RF-17.03} & Aprobación de dos perfiles para cambios de parámetros con impacto operacional & Administrador del sistema & El sistema debe requerir aprobación de un segundo perfil para todo cambio de parámetro con impacto operacional, quedando registrado con su justificación. & Alta & RT-16.03 \\
\textbf{RF-17.04} & Registro de auditoría inalterable & Administrador del sistema & El sistema debe registrar toda operación que cree, modifique o elimine información, con identificación de la persona o sistema que la ejecutó, fecha/hora con zona horaria, origen, valores anteriores y posteriores. El registro de auditoría debe ser inalterable y no modificable por ningún perfil. & Muy Alta & RT-16.06, RT-16.07 \\
\textbf{RF-17.05} & Consulta y exportación de auditoría desde la interfaz & Administrador del sistema / Auditor & El CLIENTE debe poder consultar y exportar la auditoría desde la interfaz, con filtros por persona, período, entidad y tipo de operación, sin requerir acceso a la base de datos. & Muy Alta & RT-16.08 \\
\textbf{RF-17.06} & Soporte de flujos de trabajo configurables & Administrador del sistema & El sistema debe soportar flujos de trabajo con estados, transiciones, responsables, plazos, escalamiento automático por vencimiento y delegación por ausencia, configurables por el CLIENTE sin desarrollo (al menos en responsables, plazos y niveles de aprobación). & Alta & RT-16.11, RT-16.12 \\
\textbf{RF-17.07} & Bandeja de tareas unificada & Persona usuaria (interna) & El sistema debe proveer una bandeja de tareas unificada para cada persona usuaria, donde se visualicen todas sus tareas pendientes (aprobaciones, revisiones, excepciones) con priorización y alerta de vencimiento. & Alta & RT-16.13 \\
\textbf{RF-17.08} & Gestión documental con versionado y control de acceso & Administrador del sistema & El sistema debe gestionar documentos con versionado, metadatos, control de acceso, búsqueda por contenido y por metadato, y previsualización sin descarga. Los documentos deben almacenarse cifrados con verificación de integridad. & Alta & RT-16.15, RT-16.16 \\
\textbf{RF-17.09} & Generación de documentos a partir de plantillas & Administrador del sistema & El sistema debe generar documentos a partir de plantillas administrables por el CLIENTE, con datos de la transacción y salida en formato abierto PDF, DOCX, EXCELL. & Media & RT-16.19 \\
\textbf{RF-17.10} & Notificaciones multicanal (correo, SMS, notificación en app) & Administrador del sistema & El sistema debe enviar notificaciones por al menos tres canales: correo electrónico, notificación en la aplicación y SMS o mensajería instantánea, con plantillas administrables por el CLIENTE. & Alta & RT-16.20, RT-16.21 \\
\textbf{RF-17.11} & Búsqueda global con indexación de texto completo & Persona usuaria (interna) & El sistema debe disponer de búsqueda global con indexación de texto completo con tolerancia a errores de escritura, filtros facetados y respeto del control de acceso de la persona que busca. & Alta & RT-16.27 \\
\textbf{RF-17.12} & Listados ordenables, filtrables, paginados y exportables & Persona usuaria (interna) & El sistema debe presentar listados ordenables, filtrables, paginados y exportables en formatos abiertos, con el filtro aplicado reflejado en la exportación. & Alta & RT-16.28 \\
\textbf{RF-17.13} & Portal público con catálogo y contacto (sin autenticación) & Usuario no autenticado & El sistema debe disponer de un portal público con la información que el caso determine (catálogo sin precios/stock, contacto, información institucional), accesible sin autenticación, con estándares de accesibilidad y desempeño. & Media & RT-16.31 \\
\textbf{RF-18.01} & Soporte de firma electrónica avanzada (Ley 19.799) & Administrador del sistema / Persona usuaria & El sistema debe soportar firma electrónica conforme a la Ley N\ensuremath{^\circ} 19.799, con firma avanzada para los actos que el caso lo requiera (guía de despacho electrónica, acuse de recibo, contratos), y verificación de validez del certificado al momento de la firma, generando sello de tiempo y conservando la evidencia de firma. & Muy Alta & RT-16.17, RT-16.18 \\
\textbf{RF-18.02} & Configuración de preferencias de notificación por usuario & Persona usuaria (interna o externa) & Cada persona usuaria debe poder configurar sus preferencias de canal (correo, SMS, notificación en app) y de frecuencia para las notificaciones, respetando las que el CLIENTE defina como obligatorias. & Media & RT-16.22 \\
\textbf{RF-18.03} & Envío asíncrono de notificaciones con reintento y registro & Administrador del sistema & El sistema debe enviar notificaciones de forma asíncrona, con reintento ante falla, control de duplicados y registro de entrega, apertura y error por cada mensaje. & Alta & RT-16.23 \\
\textbf{RF-13.04} & Modelo de dominio de negocio & Equipo de proyecto & El PROPONENTE debe presentar un modelo de dominio del negocio del caso, con las entidades principales, sus relaciones y los eventos de negocio que las modifican. & Uso de la expresión ambigua "o similar" al referirse a la notación del modelo de dominio (viola la regla de evitar frases que dejan abierta la interpretación). No es verificable qué otra notación sería aceptada. Se recomienda: "...entregado con notación UML". &  \\
\textbf{RF-17.09} & Generación de documentos a partir de plantillas & Administrador del sistema & El sistema debe generar documentos a partir de plantillas administrables por el CLIENTE, con datos de la transacción y salida en formato abierto (PDF, DOCX, etc.). & Contiene "etc." en "formato abierto (PDF, DOCX, etc.)", lo que viola explícitamente la regla de no usar "etc." en la redacción de requerimientos (deja el alcance indefinido). Se recomienda enumerar taxativamente los formatos aceptados. &  \\
\textbf{RF-17.11} & Búsqueda global con indexación de texto completo & Persona usuaria (interna) & El sistema debe disponer de búsqueda global con indexación de texto completo, tolerancia a errores de escritura, filtros facetados y respeto del control de acceso de la persona que busca. & El resultado esperado usa el adjetivo ambiguo "búsqueda rápida" sin métrica asociada, lo que lo hace no verificable (viola la regla de evitar adjetivos ambiguos tipo "rápido"). Se recomienda fijar un umbral, ej. "tiempo de respuesta $\le$ 2 segundos para el percentil 95". &  \\
\end{longtable}



\subsubsection{Requerimientos no funcionales de bases}
\begin{longtable}{p{1.7cm} p{2.3cm} p{1.9cm} p{5.0cm} p{1.1cm} p{1.9cm}}
\caption{Requerimientos no funcionales de Bases (RT)}\label{tab:bt_rnf}\\
\toprule
\textbf{ID} & \textbf{T\'itulo} & \textbf{Actor} & \textbf{Descripci\'on} & \textbf{Prioridad} & \textbf{Origen} \\
\midrule
\endfirsthead
\toprule
\textbf{ID} & \textbf{T\'itulo} & \textbf{Actor} & \textbf{Descripci\'on} & \textbf{Prioridad} & \textbf{Origen} \\
\midrule
\endhead
\bottomrule
\endlastfoot
\textbf{RNF-13.01} & Disponibilidad del sitio on-premise durante corte de enlace (24h) & — & El componente on-premise debe operar de forma autónoma y degradada ante la pérdida total del enlace con la nube, durante un período mínimo de 24 horas continuas, continuando registrando las transacciones operacionales críticas de forma local, con integridad garantizada y sin pérdida de datos. & Muy Alta & RT-03.10, RT-03.11, RT-03.12, RT-03.13 \\
\textbf{RNF-13.02} & Declaración de funciones no disponibles en modo desconectado & Equipo de proyecto & El PROPONENTE debe declarar qué funciones NO estarán disponibles en modo desconectado y qué procedimiento manual las suple. & Muy Alta & RT-03.13 \\
\textbf{RNF-13.03} & Redundancia de equipos on-premise críticos &  & Los equipos on-premise críticos deben ser redundantes. & Muy Alta & RT-03.14 \\
\textbf{RNF-13.04} &  &  & El almacenamiento local debe tolerar la falla de al menos un disco & Muy Alta & RT-03.14 \\
\textbf{RNF-13.05} &  &  & el PROPONENTE debe declarar el nivel RAID escogido y justificarlo frente a las alternativas. & Muy Alta & RT-03.14 \\
\textbf{RNF-13.06} & Endurecimiento de sistemas on-premise según CIS Benchmarks & Administrador del sistema & Los sistemas on-premise deben endurecerse conforme a los CIS Benchmarks aplicables, con gestión centralizada de parches y ventana de aplicación acordada con el CLIENTE. & Alta & RT-03.15 \\
\textbf{RNF-13.07} & Enlace redundante entre on-premise y nube & — & El enlace entre el sitio on-premise y la nube debe ser redundante, con caminos físicos y proveedores distintos, y conmutación automática con tiempo de conmutación declarado ($\le$ 5 min). & Muy Alta & RT-03.17 \\
\textbf{RNF-13.08} & Dimensionamiento de ancho de banda por sitio & Equipo de proyecto & El PROPONENTE debe dimensionar el ancho de banda requerido por sitio, en régimen normal y en peak, y justificarlo con el cálculo de volumen de transacciones y de datos. & Alta & RT-03.20 \\
\textbf{RNF-13.09} & Estudio de cobertura de red inalámbrica en sitios operacionales (incluye cámaras de frío) & Equipo de proyecto & El PROPONENTE debe realizar un estudio de cobertura de red inalámbrica en los centros de distribución, que incluya expresamente el interior de las cámaras de refrigerado y de congelado, hoy sin señal, y proponer la solución técnica para resolverlo. & Muy Alta & RT-03.23 (Según caso) \\
\textbf{RNF-14.01} & Cifrado en tránsito con TLS 1.3 & — & Todo el tráfico en tránsito debe emplear TLS 1.3, con prohibición expresa de TLS 1.0 y 1.1, conjuntos de cifrado modernos, HSTS con precarga y gestión automatizada de certificados con rotación y alerta anticipada de vencimiento. & Muy Alta & RT-11.08 \\
\textbf{RNF-14.02} & Cifrado en reposo de todos los datos & — & La totalidad de los datos en reposo debe estar cifrada, con claves gestionadas en un servicio de gestión de claves o en un módulo de seguridad de hardware, política de rotación declarada y separación de funciones en la custodia de claves. & Muy Alta & RT-11.09 \\
\textbf{RNF-14.03} & Segregación de red en la nube (subredes privadas) & — & En la nube, la red debe segmentarse por capas, con subredes privadas para aplicación y datos, y exposición pública restringida a la capa de borde. Ningún componente de datos debe ser alcanzable desde Internet. & Muy Alta & RT-03.04 \\
\textbf{RNF-14.04} & SIEM con casos de uso de detección específicos del negocio & Administrador de seguridad & Los eventos de seguridad deben registrarse de forma centralizada e inalterable, y correlacionarse en una plataforma SIEM con casos de uso de detección definidos específicamente para el proceso de negocio del caso (ej. intentos de acceso a trazabilidad de lote, modificaciones de inventario no autorizadas). & Alta & RT-11.14, RT-11.15 \\
\textbf{RNF-14.05} & Detección y respuesta en puntos finales (EDR) & Administrador de seguridad & Se debe implementar detección y respuesta en puntos finales (EDR) y en cargas de trabajo, tanto en nube como on-premise, con cobertura 24$\times$7. & Alta & RT-11.16 \\
\textbf{RNF-14.06} & Pruebas de intrusión anuales por tercero independiente & Administrador de seguridad & Se deben ejecutar pruebas de intrusión (pentesting) por un tercero independiente del ADJUDICATARIO, anualmente y antes de cada paso a producción, con entrega íntegra del informe al CLIENTE y plan de remediación con plazos. & Muy Alta & RT-11.20 \\
\textbf{RNF-14.07} & Análisis estático de código, SCA y escaneo de contenedores en CI/CD & Equipo de desarrollo & El flujo de integración continua debe incorporar análisis estático de código (SAST), análisis de composición de software (SCA), análisis dinámico (DAST) y escaneo de imágenes de contenedor, con criterios de bloqueo automático del despliegue ante hallazgos críticos. & Muy Alta & RT-11.22 \\
\textbf{RNF-14.08} & Inventario de componentes de software (SBOM) en formato CycloneDX/SPDX & Equipo de desarrollo & Cada versión liberada debe acompañarse de su inventario de componentes de software en formato CycloneDX o SPDX, entregado al CLIENTE. & Alta & RT-11.23 \\
\textbf{RNF-14.09} & Verificación de procedencia de artefactos según SLSA nivel 3 & Equipo de desarrollo & Los artefactos deben firmarse y su procedencia debe verificarse conforme a SLSA nivel 3 o superior. & Alta & RT-11.24 \\
\textbf{RNF-14.10} & Prohibición de datos productivos en ambientes no productivos & Equipo de proyecto & Queda prohibido el uso de datos productivos reales en ambientes no productivos sin anonimización o seudonimización verificable. & Muy Alta & RT-11.25 \\
\textbf{RNF-15.01} & Política de sesión: duración, inactividad, renovación, revocación & — & El sistema debe implementar política de sesión con duración máxima, caducidad por inactividad, renovación de la credencial de sesión tras la autenticación, revocación inmediata y control de sesiones concurrentes. & Muy Alta & RT-12.07 \\
\textbf{RNF-15.02} & Credenciales de sesión firmadas y de vida breve & — & Las credenciales de sesión deben ser firmadas y de vida breve, con credencial de refresco rotatoria. Queda prohibido transportar identificadores de sesión en la ruta de la dirección web (URL). & Muy Alta & RT-12.08 \\
\textbf{RNF-15.03} & Registro completo del ciclo de vida de la identidad (auditoría) & — & El sistema debe registrar auditoría completa del ciclo de vida de la identidad: creación, modificación, elevación, bloqueo y baja de cuentas, con no repudio y retención declarada. & Muy Alta & RT-12.09 \\
\textbf{RNF-16.01} & Indicadores de nivel de servicio medidos sobre experiencia real del usuario & Equipo de operaciones & Los indicadores de nivel de servicio deben medirse sobre la experiencia real de la persona usuaria (percentil 95) y no sobre pruebas sintéticas, sin perjuicio de que estas se empleen como complemento. & Muy Alta & RT-14.03 \\
\textbf{RNF-16.02} & Libro de operación y guía de resolución documentados & Equipo de operaciones & El PROPONENTE debe documentar un libro de operación y una guía de resolución para cada escenario de falla previsible, con automatización progresiva de las tareas repetitivas. & Alta & RT-14.05 \\
\textbf{RNF-16.03} & Registros sin datos personales sensibles ni credenciales & — & Los registros (logs) de la solución no deben contener datos personales sensibles ni credenciales, y su acceso debe estar controlado y auditado. & Muy Alta & RT-14.07 \\
\textbf{RNF-17.01} & Retención de auditoría mínima de 5 años & — & El período de retención de la auditoría debe ser el que fije el caso y, en su defecto, no inferior a cinco años. & Muy Alta & RT-16.10 \\
\textbf{RNF-17.02} & Exportaciones de gran volumen procesadas de forma asíncrona & — & Las exportaciones de gran volumen deben procesarse de forma asíncrona, con notificación al completarse y sin bloquear la sesión del usuario. & Alta & RT-16.29 \\
\textbf{RNF-17.03} & Portal público resistente a picos de tráfico & — & El portal público debe resistir picos de tráfico sin degradar los servicios transaccionales internos, mediante aislamiento de recursos y caché. & Alta & RT-16.34 \\
\textbf{RNF-19.01} & Cálculo de capacidad con supuestos de usuarios concurrentes y TPS & Equipo de proyecto & El PROPONENTE debe presentar el cálculo de capacidad que sustenta su dimensionamiento, con los supuestos de usuarios concurrentes, transacciones por segundo, volumen de datos y crecimiento anual, tomados de la volumetría del caso. & Muy Alta & RT-09.01 \\
\textbf{RNF-19.02} & Escalamiento horizontal automático de capas de aplicación e integración & — & Las capas de aplicación e integración deben escalar horizontalmente de forma automática, con umbrales, límites superiores y costo asociado declarados en la oferta. & Muy Alta & RT-02.10, RT-09.04 \\
\textbf{RNF-19.03} & Identificación del cuello de botella al crecer la carga & Equipo de proyecto & El PROPONENTE debe identificar el componente que primero se convertirá en cuello de botella al crecer la carga y explicar cómo lo detectará y cómo lo resolverá. & Alta & RT-09.05 \\
\textbf{RNF-19.04} & Pruebas de carga en Preproducción con 1.5$\times$ peak y pruebas de estrés & Equipo de proyecto & Se deben ejecutar pruebas de carga sobre Preproducción con un volumen equivalente a 1,5 veces el peak declarado, y pruebas de estrés hasta identificar el punto de quiebre de la solución. & Muy Alta & RT-09.06, RT-09.07 \\
\textbf{RNF-20.01} & Disponibilidad mensual mínima del 99.9\% para servicios críticos & — & La solución debe alcanzar una disponibilidad mensual mínima de 99.9\% para los servicios clasificados como críticos, medida sobre la transacción de negocio de extremo a extremo. & Muy Alta & RT-10.01 \\
\textbf{RNF-20.02} & Clasificación de servicios por criticidad (crítico, alto, medio, bajo) & Equipo de proyecto & El PROPONENTE debe clasificar cada servicio de la solución en crítico, alto, medio o bajo, justificando la clasificación con el impacto operacional de su indisponibilidad, y aplicar a cada uno el nivel de servicio correspondiente. & Muy Alta & RT-10.02 \\
\textbf{RNF-20.03} & Plan de continuidad del negocio conforme a ISO 22301 & Equipo de proyecto & El PROPONENTE debe elaborar un plan de continuidad del negocio conforme a ISO 22301, con análisis de impacto en el negocio, escenarios de contingencia, procedimientos manuales de respaldo y criterios de activación. & Muy Alta & RT-10.03 \\
\textbf{RNF-20.04} & Pruebas de resiliencia con inyección de fallas & Equipo de proyecto & Se deben ejecutar pruebas de resiliencia mediante inyección controlada de fallas (caída de instancia, de zona, de dependencia externa, latencia elevada, saturación de disco) antes de cada paso a producción y al menos una vez por semestre durante la Operación. & Alta & RT-10.07 \\
\textbf{RNF-20.05} & Documentación de comportamiento ante falla de dependencias externas & Equipo de proyecto & El PROPONENTE debe documentar, por cada dependencia externa, el comportamiento de la solución cuando esa dependencia no responde, responde con error o responde con lentitud. & Alta & RT-10.08 \\
\textbf{RNF-20.06} & Site secundario para DR: RTO $\le$ 4h, RPO $\le$ 15 min & — & El PROPONENTE debe habilitar un sitio secundario para recuperación ante desastres, con RTO $\le$ 4 horas y RPO $\le$ 15 minutos para los servicios críticos, y probar la conmutación al menos 2 veces al año. & Muy Alta & RT-07.04, RT-07.07 \\
\textbf{RNF-20.07} & Política de respaldo 3-2-1-1-0 (3 copias, 2 medios, 1 fuera de sitio, 1 inmutable, 0 errores) & — & La política de respaldo debe seguir el esquema 3-2-1-1-0: tres copias, en dos medios distintos, una fuera de sitio, una inmutable o fuera de línea y cero errores de verificación de restauración. & Muy Alta & RT-07.09, RT-07.12 \\
\textbf{RNF-21.01} & Centro de operaciones de red (NOC) 24$\times$7$\times$365 & — & El ADJUDICATARIO debe disponer de un centro de operaciones de red (NOC) con cobertura 24$\times$7$\times$365, propio o subcontratado, y declarar su ubicación, dotación por turno y procedimientos. & Muy Alta & RT-21.01 \\
\textbf{RNF-21.02} & Gerente de servicio dedicado para el CLIENTE & — & El ADJUDICATARIO debe designar un gerente de servicio dedicado como contraparte permanente del CLIENTE durante la fase de Operación. & Muy Alta & RT-21.02 \\
\textbf{RNF-21.03} & Centro de atención telefónica: 80\% en 20 segundos, 70\% resolución primer contacto & — & El centro de atención debe cumplir un tiempo de atención al usuario final de 80\% antes de 20 segundos y una resolución al primer contacto de al menos 70\%. La tasa de abandono no debe superar el 5\%. & Muy Alta & RT-21.06 \\
\textbf{RNF-21.04} & Horario de atención ampliado a 24$\times$7 para incidentes críticos y ventana de despacho (05:30-07:00) & — & El horario mínimo de atención será de 08:00 a 20:00 en días hábiles, ampliado a 24$\times$7 para los incidentes de severidad crítica y para la ventana operacional crítica (despacho de 05:30 a 07:00). & Muy Alta & RT-21.07 (Según caso) \\
\textbf{RNF-21.05} & Dimensionamiento de la dotación del centro de atención con teoría de colas o Erlang C & Equipo de proyecto & El PROPONENTE debe dimensionar la dotación del centro de atención con fundamento cuantitativo, empleando teoría de colas o el modelo Erlang C, a partir del volumen de contactos proyectado del caso. & Alta & RT-21.14 \\
\textbf{RNF-21.06} & Canal único de registro de incidentes con ticket y seguimiento & — & Debe existir un canal único de registro de incidentes y solicitudes, con número de ticket, clasificación por severidad y seguimiento del ciclo de vida completo hasta el cierre conforme. & Alta & RT-21.15 \\
\textbf{RNF-21.08} & Bolsa anual de horas de mantención evolutiva & — & El PROPONENTE debe declarar el tamaño de la bolsa anual de horas de mantención evolutiva, su composición por perfil y el procedimiento de solicitud, estimación, aprobación y liquidación. & Alta & RT-21.19 \\
\textbf{RNF-21.09} & Planificación anual de actualizaciones de versiones de componentes & — & Las actualizaciones de versión de los componentes de base (sistemas operativos, bases de datos, frameworks) deben planificarse anualmente, con ventana acordada y plan de reversión. & Alta & RT-21.22 \\
\textbf{RNF-22.01} & Plan de capacitación estructurado por perfil (usuarios, administradores, técnicos) & Equipo de proyecto & El PROPONENTE debe presentar un plan de capacitación estructurado por perfil: personas usuarias finales, usuarias avanzadas, administradoras, equipo técnico y soporte de niveles 1 y 2. & Muy Alta & RT-22.01 \\
\textbf{RNF-22.03} & Capacitación sin afectar la operación: programación por turnos y horarios acordados & — & La capacitación no podrá afectar la operación del CLIENTE: se programará por turnos y en horarios acordados, considerando la estacionalidad (peaks de septiembre y diciembre) y la ventana operacional (despacho de mañana, preparación nocturna). & Muy Alta & RT-22.04 (Según caso) \\
\textbf{RNF-22.04} & Evaluación y certificación de usuarios administradores y equipo técnico & Equipo de proyecto & Las personas usuarias administradoras y el equipo técnico del CLIENTE deben ser evaluados y certificados como condición para el cierre de cada marcha blanca. & Alta & RT-22.05 \\
\textbf{RNF-22.05} & Capacitación de personal nuevo del CLIENTE sin costo adicional durante la Operación & — & Durante la Operación, el ADJUDICATARIO debe ejecutar al menos dos jornadas anuales de actualización y capacitar al personal nuevo del CLIENTE sin costo adicional. & Alta & RT-22.06 \\
\textbf{RNF-22.06} & Acompañamiento y mentoría al equipo técnico del CLIENTE durante 6 meses post producción & — & El ADJUDICATARIO debe proveer acompañamiento y mentoría al equipo técnico del CLIENTE durante al menos seis meses posteriores al paso a producción de la Etapa 2. & Muy Alta & RT-22.07 \\
\textbf{RNF-23.01} & Especificación de hardware de terreno con marca, modelo, cantidad y características & Equipo de proyecto & El PROPONENTE debe especificar cada dispositivo de terreno (lectores de código, impresoras térmicas, sensores de temperatura, terminales móviles) con marca, modelo de referencia, cantidad, características mínimas, accesorios, consumibles y costo unitario estimado. & Muy Alta & RT-08.10, RT-08.11 \\
\textbf{RNF-23.02} & Dispositivos con grado de protección IP y resistencia a caídas coherentes con el entorno & — & Los dispositivos de terreno deben declarar su grado de protección contra polvo y agua (IP) y su resistencia a caídas, coherentes con el entorno de operación (intemperie, humedad, polvo, vibración, temperatura). & Muy Alta & RT-08.12 \\
\textbf{RNF-23.03} & Ciclo de vida esperado de cada dispositivo y plan de reposición durante 56 meses & Equipo de proyecto & El PROPONENTE debe indicar el ciclo de vida esperado de cada dispositivo de terreno, la disponibilidad de repuestos y el plan de reposición durante los 56 meses del Contrato. & Alta & RT-08.13 \\
\textbf{RNF-23.04} & Estaciones de trabajo para operación y administración con monitores duales & Equipo de proyecto & El PROPONENTE debe especificar las estaciones de trabajo requeridas para la operación y administración de la plataforma, en la cantidad que determine el dimensionamiento del caso, con monitores duales, cumpliendo condiciones ergonómicas NCh 2527 y gestión centralizada con cifrado de disco. & Alta & RT-08.07, RT-08.08, RT-08.09 \\
\textbf{RNF-21.07} & Traslado a sitios alejados (Curicó, Chillán, Los Ángeles) para resolución de incidentes & — & Los especialistas de niveles 2 y 3 deben trasladarse a los sitios alejados (Curicó, Chillán, Los Ángeles) cuando la resolución lo requiera, por el medio más rápido disponible, sin alterar la atención normal del resto de los sitios. El costo del traslado está incluido en la oferta. & Usa el superlativo ambiguo "por el medio más rápido disponible", lo que viola la regla de evitar superlativos como "mejor"/"más" sin criterio objetivo de comparación. Se recomienda especificar el medio de transporte o un tiempo máximo de traslado. &  \\
\textbf{RNF-22.02} & Material de capacitación en español, editable y de propiedad del CLIENTE & — & Todo el material de capacitación debe entregarse en español, en formato editable (Word, PowerPoint, etc.) y de propiedad del CLIENTE: manuales por perfil, guías rápidas, preguntas frecuentes, videos tutoriales y base de conocimiento consultable. & Contiene "etc." en "formato editable (Word, PowerPoint, etc.)", lo que viola explícitamente la regla de no usar "etc.". Se recomienda enumerar los formatos aceptados de forma taxativa. &  \\
\end{longtable}



